\chapter{Introduction}

Modern analyses of deep networks mix high-dimensional probability, gradient-flow ODEs, and parameterization-dependent scaling limits. Informal papers routinely hide constants, interchange limits, and leave feature-learning claims un-checked.

OpenTheory is a dual-track project in the sense of Massot's Lean blueprints~\cite{tao2023pfr}: every claim below is written for humans in \LaTeX{} and, when a Lean name is given, mapped to a declaration in the \texttt{OpenTheory} library. Border colour on the dependency graph records whether the \emph{statement} is in Lean (\texttt{\textbackslash leanok} on the environment). Fill colour records whether the \emph{proof} is kernel-checked (\texttt{\textbackslash leanok} inside \texttt{proof}).

The chapters follow the dependency order of the library. Data states the hypotheses. Architectures define the network and what it can represent. Parameterization fixes how width enters initialization and learning rates. Representation defines features, the data Gram matrix, and alignment. Optimization is the training dynamics of the loss. Generalization passes from empirical risk to population risk. Dynamics describes how representations move.

The only axioms permitted besides Mathlib's \texttt{propext}, \texttt{Classical.choice}, and \texttt{Quot.sound} are \texttt{sorryAx} on explicitly open targets. Custom axioms fail CI.

\chapter{Data}

\begin{definition}[Labeled sample]\label{def:sample}
\lean{OpenTheory.Sample}
\leanok
A sample of size $n\ge 1$ in dimension $d$ is a family $(x_i,y_i)_{i=1}^n$ with
\[
\sum_{j=1}^d x_{i,j}^2 \le 1, \qquad |y_i|\le 1
\]
for every $i$.
\end{definition}

\begin{definition}[Spectral gap]\label{def:spectral_gap}
\lean{OpenTheory.SpectralGap}
\leanok
A spectral gap is a triple $(\lambda_{\mathrm{lead}},\lambda_{\mathrm{bulk}},\gamma)$ with $\gamma>0$, $\lambda_{\mathrm{bulk}}\ge 0$, and
\[
\lambda_{\mathrm{lead}} \ge \lambda_{\mathrm{bulk}} + \gamma.
\]
\end{definition}

\begin{definition}[Spiked covariance]\label{def:spiked_model}
\lean{OpenTheory.SpikedModel}
\leanok
\uses{def:spectral_gap}
A spiked model in dimension $d$ records a signal rank $k$ with $1\le k\le d$, a signal strength, a noise level, and a spectral gap. The signal coordinates are the first $k$ coordinates of the input.
\end{definition}

\chapter{Architectures}

\begin{definition}[Activation function]\label{def:activation}
\lean{OpenTheory.ActivationFunction}
\leanok
A map $\sigma:\R\to\R$ is an admissible activation if it is differentiable and Lipschitz: there exists $L\in\R$ such that
\[
|\sigma(x)-\sigma(y)| \le L\,|x-y|
\quad\text{for all }x,y\in\R.
\]
\end{definition}

\begin{definition}[Two-layer network]\label{def:two_layer_net}
\lean{OpenTheory.TwoLayerNet, OpenTheory.TwoLayerNet.forward}
\leanok
\uses{def:activation}
For input dimension $d$ and width $m$, a two-layer network is a triple $(W_1,W_2,\sigma)$ with $W_1\in\R^{m\times d}$, $W_2\in\R^{1\times m}$, and $\sigma$ an admissible activation. The forward map is
\[
f_\theta(x)
=
\sum_{k=1}^m W_{2,k}\,
\sigma\Bigl(\sum_{j=1}^d W_{1,kj}\,x_j\Bigr).
\]
\end{definition}

\begin{definition}[Residual step]\label{def:residual}
\lean{OpenTheory.ResidualBranch, OpenTheory.residualStep}
\leanok
\uses{def:activation}
A residual branch is a square matrix $W$ together with an admissible activation. The step of scale $\alpha\in\R$ is
\[
h \mapsto h + \alpha\,\sigma(Wh).
\]
\end{definition}

\begin{definition}[Empirical approximation error]\label{def:approx_error}
\lean{OpenTheory.coordinateLift, OpenTheory.NonPolynomial, OpenTheory.empiricalApproxError}
\leanok
\uses{def:two_layer_net, def:sample}
For a network $f_\theta$, a target $g$, and a sample $(x_i)$, the empirical approximation error is
\[
\frac1n\sum_{i=1}^n \bigl|f_\theta(x_i)-g(x_i)\bigr|.
\]
$\mathtt{coordinateLift}$ reads the first $k$ coordinates of a vector in $\R^d$. $\mathtt{NonPolynomial}$ says $\sigma$ agrees with no polynomial.
\end{definition}

\begin{lemma}[Pointwise control of approximation error]\label{lem:approx_error_pointwise}
\lean{OpenTheory.empiricalApproxError_le_of_pointwise}
\leanok
\uses{def:approx_error}
If $|f_\theta(x_i)-g(x_i)|\le\varepsilon$ for every sample index $i$, then the empirical approximation error is at most $\varepsilon$.
\end{lemma}
\begin{proof}
\leanok
The average of numbers each at most $\varepsilon$ is at most $\varepsilon$. Multiply the resulting sum bound by $n^{-1}$.
\end{proof}

\begin{lemma}[Finite-sample expressivity on the spike]\label{lem:expressivity_subspace}
\lean{OpenTheory.forward_zero_of_act_zero, OpenTheory.biasless_zero_sample_obstruction, OpenTheory.subspaceLiftNet, OpenTheory.subspaceLiftNet_forward_eq, OpenTheory.empiricalApproxError_subspaceLiftNet, OpenTheory.exists_exact_sample_interpolator_of_anchor, OpenTheory.exists_subspace_approximation_of_signal_net}
\leanok
\uses{def:spiked_model, def:approx_error, def:activation}
Let the target be a function of the first $k = \mathrm{signalRank}$ coordinates of a spiked model. Embedding a signal-subspace network $\mathrm{net}_{\mathrm{sig}}$ into $\R^d$ via \texttt{subspaceLiftNet} (zeroing the $d-k$ bulk coordinates of $W_1$) satisfies $(\mathrm{subspaceLiftNet}\,\mathrm{net}_{\mathrm{sig}})(x) = \mathrm{net}_{\mathrm{sig}}(x_{1:k})$ for all $x\in\R^d$, preserving empirical approximation error exactly while eliminating all bulk-coordinate sensitivity. Moreover, when dual separating directions exist on the signal coordinates with $\sigma(0)=0$ and $\sigma(c)\ne 0$, a width-$n$ lifted network achieves exact zero empirical error (\texttt{exists\_exact\_sample\_interpolator\_of\_anchor}); conversely, for biasless networks with $\sigma(0)=0$ on a zero sample, we record the exact obstruction $\mathrm{empiricalApproxError} = |\mathrm{target}(0)|$ (\texttt{biasless\_zero\_sample\_obstruction}).
\end{lemma}
\begin{proof}
\leanok
Split the sum over $\mathrm{Fin}\,d$ into the first $k$ signal coordinates and the remaining $d-k$ zeroed coordinates to establish \texttt{subspaceLiftWeights\_dot\_eq} and \texttt{subspaceLiftNet\_forward\_eq}. For exact sample interpolation, set $W_{2,0,i} = \mathrm{target}(z_i)/\sigma(c)$ so each neuron activates solely on sample $i$.
\end{proof}

\chapter{Parameterization}

\begin{definition}[Maximal Update Parametrization]\label{def:mup_scaling}
\lean{OpenTheory.MuPScaling, OpenTheory.MuPScaling.standard, OpenTheory.MuPScaling.input_dim_pos, OpenTheory.MuPScaling.width_pos}
\leanok
\uses{def:two_layer_net}
Under $\mup$~\cite{yang2021tensor4,yang2022tensor5}, the default coordinate scaling is
\[
W_1\sim\mathcal{N}(0,1/d),\qquad
W_2\sim\mathcal{N}(0,1/m),
\]
with learning rates $\eta_{W_1}=\Theta(1)$ and $\eta_{W_2}=\Theta(1/m)$. The structure stores those four scale factors together with $d>0$ and $m>0$. \texttt{MuPScaling.standard} installs this choice.
\end{definition}

\begin{definition}[Neural tangent parameterization]\label{def:ntk_scaling}
\lean{OpenTheory.NTKScaling, OpenTheory.NTKScaling.standard, OpenTheory.NTKScaling.input_dim_pos, OpenTheory.NTKScaling.width_pos}
\leanok
\uses{def:two_layer_net}
In the neural tangent parameterization~\cite{jacot2018neural,yang2021tensor4}, \texttt{NTKScaling.standard} sets both initialization standard deviations to $1/\sqrt{\mathrm{fan\_in}}$ and both learning rates to $1$. At infinite width the hidden features stay at initialization. The feature-learning conjecture is the statement that $\mup$ leaves this regime.
\end{definition}

\chapter{Representation}

\begin{definition}[Hidden features]\label{def:features}
\lean{OpenTheory.hiddenFeatures, OpenTheory.featureMatrix, OpenTheory.forward_eq_readout}
\leanok
\uses{def:two_layer_net, def:sample}
The hidden representation of a two-layer net on $x\in\R^d$ is the vector $\sigma(W_1 x)\in\R^m$. On a sample, these vectors are the rows of the feature matrix. The network output is the readout of that vector by $W_2$.
\end{definition}
\begin{proof}
\leanok
The readout identity is the definition of \texttt{forward}, unfolded.
\end{proof}

\begin{definition}[Data Gram matrix]\label{def:gram}
\lean{OpenTheory.dataGram}
\leanok
\uses{def:sample}
The empirical Gram matrix of a sample is $G\in\R^{d\times d}$ with
\[
G_{ij} = \sum_{s=1}^n x_{s,i}\, x_{s,j}.
\]
\end{definition}

\begin{definition}[Alignment]\label{def:alignment}
\lean{OpenTheory.frobeniusSq, OpenTheory.IsOrthoProjector, OpenTheory.alignment}
\leanok
\uses{def:features, def:gram}
Let $H$ be a feature matrix with rows indexed by the sample and let $P$ be a matrix on feature coordinates. Alignment is the fraction of Frobenius energy that survives $P$:
\[
a(H,P) = \frac{\lVert HP\rVert_F^2}{\lVert H\rVert_F^2},
\]
and $a(H,P)=0$ when $H=0$. An orthogonal projector satisfies $P^2=P$ and $P^\top=P$. The projector used in later chapters is the orthogonal projector onto a leading eigenspace of the Gram matrix; that identification is not yet a definition in Lean.
\end{definition}

\begin{lemma}[Alignment is non-negative]\label{lem:alignment_nonneg}
\lean{OpenTheory.alignment_nonneg}
\leanok
\uses{def:alignment}
For every feature matrix $H$ and every matrix $P$ of the right shape, $a(H,P)\ge 0$.
\end{lemma}
\begin{proof}
\leanok
Both the numerator and the denominator are sums of squares. The zero-matrix case returns $0$.
\end{proof}

\begin{lemma}[Alignment is at most one]\label{lem:alignment_le_one}
\lean{OpenTheory.alignment_le_one, OpenTheory.orthoProjector_row_sq_sum, OpenTheory.frobeniusSq_sub_mul_orthoProjector, OpenTheory.frobeniusSq_mul_orthoProjector_le}
\leanok
\uses{def:alignment}
If $P$ is an orthogonal projector, then $a(H,P)\le 1$.
\end{lemma}
\begin{proof}
\leanok
By symmetry $P^\top=P$ and idempotence $P^2=P$, each row $i$ satisfies $\sum_j ((HP)_{ij})^2 = \sum_j (HP)_{ij} H_{ij}$. Expanding $\lVert H - HP\rVert_F^2$ row by row yields the exact Frobenius Pythagorean identity $\lVert H - HP\rVert_F^2 = \lVert H\rVert_F^2 - \lVert HP\rVert_F^2$. Since $\lVert H - HP\rVert_F^2 \ge 0$, $\lVert HP\rVert_F^2 \le \lVert H\rVert_F^2$, so dividing by $\lVert H\rVert_F^2 > 0$ gives $a(H,P)\le 1$.
\end{proof}

\begin{lemma}[Alignment certificate]\label{lem:alignment_certificate}
\lean{OpenTheory.alignment_feature_distance}
\leanok
\uses{lem:alignment_le_one}
Let $P$ is an orthogonal projector and $\lVert H\rVert_F \ne 0$. If $a(H,P)\ge 1-\varepsilon$ for some $\varepsilon\ge 0$, then
\[
\lVert H - HP\rVert_F^2 \le \varepsilon\,\lVert H\rVert_F^2.
\]
This is the sense in which a high alignment score says the features lie close to the subspace.
\end{lemma}
\begin{proof}
\leanok
From the Frobenius Pythagorean identity $\lVert H - HP\rVert_F^2 = \lVert H\rVert_F^2 - \lVert HP\rVert_F^2 = (1 - a(H,P))\lVert H\rVert_F^2$, the lower bound $a(H,P)\ge 1-\varepsilon$ immediately yields $\lVert H - HP\rVert_F^2 \le \varepsilon \lVert H\rVert_F^2$.
\end{proof}

\chapter{Optimization}

\begin{definition}[Empirical risk]\label{def:empirical_risk}
\lean{OpenTheory.EmpiricalRisk}
\leanok
\uses{def:two_layer_net, def:sample}
The empirical risk of a two-layer net on a sample is the mean squared error
\[
\frac1n\sum_{i=1}^n \bigl(f_\theta(x_i)-y_i\bigr)^2.
\]
\end{definition}

\begin{definition}[Descent path]\label{def:descent}
\lean{OpenTheory.DescentPath}
\leanok
\uses{def:empirical_risk}
A descent path is a trajectory $t\mapsto \theta_t$ of two-layer nets whose empirical risk admits a derivative $L'(t)\le 0$ at every time. The full gradient-flow ODE, identifying $L'$ with minus a learning rate times the loss gradient, is not yet formalized.
\end{definition}

\begin{lemma}[Exponential approach under a descent rate]\label{lem:risk_exponential}
\lean{OpenTheory.risk_exponential_approach}
\leanok
\uses{def:descent, def:mup_scaling}
Let the empirical risk of a descent path be bounded below by a floor $F$, and suppose its derivative satisfies $L'(t) \le -c\,(L(t)-F)$ for some $c>0$. Then for every $t\ge 0$,
\[
L(t) \le F + e^{-ct}\bigl(L(0)-F\bigr).
\]
The µP scale is a parameter of the statement. Proving the rate $c$ from the gradient-flow ODE is open.
\end{lemma}
\begin{proof}
\leanok
Define the rescaled excess risk $g(s) = e^{cs}(L(s)-F)$. By the product and chain rules, $g'(s) = e^{cs}(c(L(s)-F) + L'(s)) \le 0$ for all $s\in\mathbb{R}$. Monotonicity of functions with nonpositive derivative (`image_le_of_deriv_right_le_deriv_boundary`) yields $g(t)\le g(0) = L(0)-F$ for $t\ge 0$, and multiplying by $e^{-ct}>0$ completes the proof.
\end{proof}

\chapter{Generalization}

\begin{lemma}[Activation Lipschitz bound]\label{lem:lipschitz_gelu}
\lean{OpenTheory.ActivationFunction.lipschitz_bound}
\leanok
\uses{def:activation}
If $\sigma$ is an admissible activation, then $\sigma$ is Lipschitz: there exists $L$ such that $|\sigma(x)-\sigma(y)|\le L\,|x-y|$ for all real $x,y$.
\end{lemma}
\begin{proof}
\leanok
Immediate from the Lipschitz field of \cref{def:activation}.
\end{proof}

\begin{lemma}[Lipschitz bound for one layer]\label{lem:layer_lipschitz}
\lean{OpenTheory.dotProduct, OpenTheory.l2Norm, OpenTheory.abs_dotProduct_le, OpenTheory.abs_dotProduct_sub_le, OpenTheory.layer_lipschitz_of_lip, OpenTheory.layer_lipschitz}
\leanok
\uses{lem:lipschitz_gelu}
For an admissible activation, a weight vector $w\in\R^d$, and inputs $x,y$,
\[
\bigl|\sigma(\langle w,x\rangle)-\sigma(\langle w,y\rangle)\bigr|
\le L\,\|w\|_2\,\|x-y\|_2,
\]
where $L$ is a Lipschitz constant of $\sigma$.
\end{lemma}
\begin{proof}
\leanok
Apply the finite-sum Cauchy--Schwarz inequality (`Finset.sum_mul_sq_le_sq_mul_sq`) to $\langle w, x\rangle - \langle w, y\rangle = \langle w, x-y\rangle$ to obtain $|\langle w,x\rangle - \langle w,y\rangle| \le \|w\|_2\|x-y\|_2$, and compose with the activation Lipschitz bound $|\sigma(a)-\sigma(b)|\le L|a-b|$.
\end{proof}

\begin{lemma}[Lipschitz bound for the network]\label{lem:forward_lipschitz}
\lean{OpenTheory.forward_lipschitz_explicit, OpenTheory.forward_lipschitz}
\leanok
\uses{lem:layer_lipschitz, def:two_layer_net}
The map $x\mapsto f_\theta(x)$ of a two-layer net is Lipschitz with explicit constant $\sum_{i=1}^m |W_{2,i}|\,L\,\|W_{1,i}\|_2$.
\end{lemma}
\begin{proof}
\leanok
Apply the triangle inequality to the sum over hidden units $i\in\mathrm{Fin}\,m$, bound each unit by $|W_{2,i}|\,L\,\|W_{1,i}\|_2\,\|x-y\|_2$ using \cref{lem:layer_lipschitz}, and factor out $\|x-y\|_2$.
\end{proof}

\begin{lemma}[Bounded empirical means]\label{lem:bounded_mean}
\lean{OpenTheory.abs_mean_le_of_abs_le_one}
\leanok
Let $n\ge 1$ and $a_1,\dots,a_n\in\R$ with $|a_i|\le 1$ for all $i$. Then
\[
\Bigl|\frac1n\sum_{i=1}^n a_i\Bigr| \le 1.
\]
\end{lemma}
\begin{proof}
\leanok
Triangle inequality for finite sums and $|a_i|\le 1$ give $\bigl|\sum a_i\bigr|\le n$. Multiply by the positive scalar $n^{-1}$.
\end{proof}

\begin{lemma}[Deterministic ghost-sample and Rademacher-average inequalities]\label{lem:rademacher_bound}
\lean{OpenTheory.rademacherAvg, OpenTheory.symmetrization_ghost_decomp, OpenTheory.symmetrization_pointwise_diff_le, OpenTheory.symmetrization_two_mul_bound, OpenTheory.linear_rademacher_avg_le, OpenTheory.twoLayer_rademacher_avg_le}
\leanok
\uses{def:two_layer_net, lem:bounded_mean, lem:layer_lipschitz}
For any two samples $X, X'$ and sign vector $\varepsilon\in\{-1,+1\}^n$, the signed ghost-sample empirical difference decomposes as $\frac{1}{n}\sum_{i=1}^n \varepsilon_i(f(X_i)-f(X'_i)) = \widehat{\mathcal{R}}_X(\varepsilon,f) + \widehat{\mathcal{R}}_{X'}(-\varepsilon,f)$ and is bounded in absolute value by $|\widehat{\mathcal{R}}_X(\varepsilon,f)| + |\widehat{\mathcal{R}}_{X'}(-\varepsilon,f)| \le 2 B$ when $|\widehat{\mathcal{R}}|\le B$~\cite{bartlett2002rademacher}. Furthermore, linear predictors satisfy $|\widehat{\mathcal{R}}_X(\varepsilon, \langle w,\cdot\rangle)| \le \|w\|_2\,\|\frac{1}{n}\sum_i \varepsilon_i X_i\|_2$, and two-layer networks with $|W_{2,k}|\le A$ satisfy $|\widehat{\mathcal{R}}_X(\varepsilon, f_{\mathrm{net}})| \le A\sum_{k=1}^m |\widehat{\mathcal{R}}_X(\varepsilon, h_k)|$.
\end{lemma}
\begin{proof}
\leanok
Split $\varepsilon_i(f(X_i)-f(X'_i)) = \varepsilon_i f(X_i) + (-\varepsilon_i)f(X'_i)$ and apply the triangle inequality to obtain the factor-$2$ ghost-sample bound. Exchange the sample and coordinate sums for linear predictors and apply Cauchy--Schwarz; expand the readout layer for two-layer networks.
\end{proof}

\begin{lemma}[Deterministic Frobenius bounds on weight matrices]\label{lem:gaussian_operator_norm}
\lean{OpenTheory.matVec_l2Norm_le_frobenius, OpenTheory.sqrt_frobeniusSq_le_of_entry_abs_le, OpenTheory.sqrt_frobeniusSq_le_of_fanIn_entry_bound}
\leanok
\uses{def:two_layer_net, def:alignment}
Every matrix $W\in\R^{m\times d}$ satisfies $\|Wx\|_2 \le \|W\|_F\|x\|_2$, with $\|W\|_F \le \sqrt{md}\,B$ under entrywise bound $|W_{ij}|\le B$ and $\|W\|_F \le \sqrt{m} \le \sqrt{m}+\sqrt{d}$ under fan-in scaling $|W_{ij}|\le 1/\sqrt{d}$. The probabilistic statement for Gaussian matrices, $\|W\|_{\mathrm{op}} \le \sqrt{m}+\sqrt{d}+t$ with probability at least $1-2e^{-t^2/2}$~\cite{vershynin2018high}, is \emph{not} formalized here.
\end{lemma}
\begin{proof}
\leanok
Bound each row inner product $(Wx)_i = \langle W_{i,\cdot}, x\rangle$ via Cauchy--Schwarz, square and sum over $i\in\mathrm{Fin}\,m$ to obtain $\|Wx\|_2 \le \|W\|_F\|x\|_2$.
\end{proof}

\chapter{Dynamics}

\begin{lemma}[Conditional spectral bias in $\mup$ scaling]\label{lem:spectral_bias_mup}
\lean{OpenTheory.spectral_bias_mup_gap_alignment_cont, OpenTheory.spectral_bias_mup_gap_alignment_disc, OpenTheory.MultiStepSpectralBias.hasDerivAt_modeAmpCont, OpenTheory.MultiStepSpectralBias.modeAmpDisc_step, OpenTheory.MultiStepSpectralBias.signal_to_bulk_ratio_cont, OpenTheory.MultiStepSpectralBias.signal_to_bulk_ratio_disc, OpenTheory.MultiStepSpectralBias.alignment_lower_bound_of_signal_bulk, OpenTheory.MultiStepSpectralBias.spectral_bias_alignment_cont_ge, OpenTheory.MultiStepSpectralBias.spectral_bias_alignment_disc_ge, OpenTheory.MultiStepSpectralBias.spiked_model_alignment_cont_ge, OpenTheory.MultiStepSpectralBias.multistep_kernel_growth_improves_ridge_mse}
\leanok
\uses{def:mup_scaling, def:alignment, lem:alignment_certificate, def:spectral_gap, lem:feature_learning_advantage}
Under $\mup$ feature-learning scaling on a data Gram matrix with spectral gap $\gamma_{\mathrm{lead}} \ge \gamma_{\mathrm{bulk}} + \Delta\gamma$ ($\Delta\gamma > 0$), eigenmodes along eigenspace $\gamma$ evolve at continuous rate $e^{\gamma t}$ and discrete rate $(1 + \eta\gamma)^T$~\cite{rahaman2019spectral}:
\begin{enumerate}
\item \textbf{Exponential signal-to-bulk amplification:} When signal energy grows as $S(t) \ge e^{2\gamma_{\mathrm{lead}} t} S_0$ (or $(1+\eta\gamma_{\mathrm{lead}})^{2T} S_0$) and bulk energy satisfies $B(t) \le e^{2\gamma_{\mathrm{bulk}} t} B_0$ (or $(1+\eta\gamma_{\mathrm{bulk}})^{2T} B_0$), the bulk-to-signal ratio decays exponentially as $B(t)/S(t) \le (B_0/S_0)e^{-2\Delta\gamma t}$ in continuous time and $B_T/S_T \le (B_0/S_0)\bigl(\frac{1+\eta\gamma_{\mathrm{bulk}}}{1+\eta\gamma_{\mathrm{lead}}}\bigr)^{2T}$ in discrete time.
\item \textbf{Subspace alignment convergence to $1$:} For any orthogonal signal projector $P$, the Frobenius Pythagorean identity yields $a(H, P) \ge \frac{S}{S+B} \ge 1 - \frac{B}{S}$, so
\[
a(H_t, P) \ge 1 - \frac{B_0}{S_0}e^{-2\Delta\gamma t},\qquad
a(H_T, P) \ge 1 - \frac{B_0}{S_0}\Bigl(\frac{1+\eta\gamma_{\mathrm{bulk}}}{1+\eta\gamma_{\mathrm{lead}}}\Bigr)^{2T},
\]
and the relative unaligned feature error $\lVert H - HP\rVert_F^2 / \lVert H\rVert_F^2$ decays at the same exponential rate.
\item \textbf{Multi-step ridge risk improvement:} Over $T\ge 1$ steps, the amplified target kernel energy $K_T \ge (1+\eta\gamma_{\mathrm{lead}})^{2T} K_0 > K_0$ strictly reduces both the squared ridge bias and the total expected ridge MSE.
\end{enumerate}
\end{lemma}
\begin{proof}
\leanok
Combine $e^{2\gamma_{\mathrm{bulk}}t} = e^{2(\gamma_{\mathrm{bulk}}-\gamma_{\mathrm{lead}})t} e^{2\gamma_{\mathrm{lead}}t} \le e^{-2\Delta\gamma t} e^{2\gamma_{\mathrm{lead}}t}$ with the signal and bulk bounds in continuous time, and compare $1+\eta\gamma_{\mathrm{lead}} \ge 1+\eta\gamma_{\mathrm{bulk}} + \eta\Delta\gamma$ in discrete time. Apply the Frobenius Pythagorean identity $\lVert H\rVert_F^2 = \lVert HP\rVert_F^2 + \lVert H-HP\rVert_F^2$ to lower-bound $a(H,P)$ by $1 - B/S$, and invoke Lemma~\ref{lem:feature_learning_advantage} for the multi-step ridge MSE improvement.
\end{proof}

\begin{lemma}[Representation contraction in residual blocks]\label{lem:representation_contraction}
\lean{OpenTheory.resnet_step_contraction_le, OpenTheory.resnet_representation_contraction_lipschitz, OpenTheory.DeepResidualScaling.branchOutput_l2Norm_le, OpenTheory.DeepResidualScaling.residualStep_l2Norm_le, OpenTheory.DeepResidualScaling.residualStep_sqrt_depth_scaling, OpenTheory.DeepResidualScaling.deepResidual_l2Norm_le, OpenTheory.DeepResidualScaling.one_add_div_pow_le_exp, OpenTheory.DeepResidualScaling.deepResidualChain_depth_independent_exp_bound, OpenTheory.DeepResidualScaling.deepResidual_norm_window}
\leanok
\uses{lem:lipschitz_gelu, def:residual, lem:layer_lipschitz}
For a residual branch $h \mapsto h + \alpha\,\sigma(Wh)$ with $L_\sigma$-Lipschitz activation and $L_\sigma\|W\|_F \le c$:
\begin{enumerate}
\item \textbf{Single-block $1/\sqrt{L}$ and $1/L$ scaling:} With $\alpha = 1/\sqrt{L}$~\cite{hayou2021stable}, Euclidean representation distances expand by at most $1 + c/\sqrt{L}$:
\[
\|h_{\ell+1}(x)-h_{\ell+1}(y)\|_2
\le
\Bigl(1+\frac{c}{\sqrt{L}}\Bigr)
\|h_\ell(x)-h_\ell(y)\|_2.
\]
\item \textbf{Depth-independent $\exp(c)$ stability across $L$ layers:} Across an $L$-layer residual chain with branch scaling $\alpha = 1/L$, the compound amplification $(1 + c/L)^L$ is bounded uniformly for all $L\ge 1$ by $e^c$:
\[
\bigl\|\mathrm{deepResidual}_L(x) - \mathrm{deepResidual}_L(y)\bigr\|_2 \le e^c\,\|x - y\|_2,
\]
and when $\sigma(0)=0$, representation norms satisfy the two-sided non-collapse window $(1 - |\alpha|c)\|h\|_2 \le \|\mathrm{residualStep}(h)\|_2$ and $\|\mathrm{deepResidual}_L(x)\|_2 \le e^c\|x\|_2$.
\end{enumerate}
\end{lemma}
\begin{proof}
\leanok
Bound each coordinate of $\sigma(Wx)-\sigma(Wy)$ via Cauchy--Schwarz and $L_\sigma$-Lipschitzness, square and sum over rows to obtain $\|\sigma(Wx)-\sigma(Wy)\|_2 \le L_\sigma\|W\|_F\|x-y\|_2$, and apply the Euclidean triangle inequality to $x + \alpha\sigma(Wx)$. Induct over $L$ layers and apply $(1+c/L)^L \le (e^{c/L})^L = e^c$.
\end{proof}

\chapter{Checked reuse foundation}

\begin{definition}[Data and probability interfaces]\label{def:reuse_data}
\lean{OpenTheory.Dataset, OpenTheory.Parameters, OpenTheory.Population, OpenTheory.IIDSample}
\leanok
Finite observations are independent of boundedness assumptions. Parameters use
Euclidean geometry. A population is a probability law; IID samples have independent
coordinates with that common law. Population-risk integrability is an explicit
separate predicate, not inferred from writing an integral.
\end{definition}

\begin{lemma}[Finite-data ridge training]\label{lem:reuse_ridge}
\lean{OpenTheory.Dataset.ridgeLoss_eq, OpenTheory.Dataset.ridge_training, OpenTheory.Ridge.admissible_step}
\leanok
\uses{def:reuse_data}
For a nonempty sample and $\lambda>0$, the objective
$\frac{1}{2n}\sum_i(\langle\theta,x_i\rangle-y_i)^2+
\frac{\lambda}{2}\|\theta\|^2$ has a minimizer. Let $A$ be its observation
operator and $L=\|A^*A/n+\lambda I\|+|\lambda|$.
For $\eta>0$ and $\eta L^2<2\lambda$, exact gradient descent satisfies
$\|\theta_k-\theta_*\|^2\leq q^k\|\theta_0-\theta_*\|^2$,
where $q=1-2\eta\lambda+\eta^2L^2<1$.
The step $\eta=\lambda/L^2$ is admissible.
\end{lemma}
\begin{proof}
\leanok
Identify the actual gradient using Mathlib calculus. Positive regularization makes
the finite-dimensional Hessian injective and hence surjective, giving a stationary
point. Quadratic expansion proves minimization. Apply TorchLean's strongly monotone
Lipschitz-operator gradient-descent theorem.
\end{proof}

\begin{lemma}[Reused neural semantics and global bounds]\label{lem:reuse_mlp}
\lean{OpenTheory.Torch.legacy_correspondence, OpenTheory.Torch.MLP.lipschitz, OpenTheory.Torch.MLP.residual_lipschitz}
\leanok
A biased two-layer ReLU MLP built from TorchLean affine layers is globally
$\|W_2\|_F\|W_1\|_F$-Lipschitz. Its residual version is
$(1+\|W_2\|_F\|W_1\|_F)$-Lipschitz. A zero-bias instantiation with the legacy
activation agrees exactly with the original forward map.
\end{lemma}
\begin{proof}
\leanok
Reuse TorchLean's Euclidean tensor metric, matrix Frobenius bound and ReLU
Lipschitz theorem; compose with Mathlib's Lipschitz lemmas. Expand coordinates
for the legacy correspondence. No differentiability of ReLU at zero is assumed.
\end{proof}

\begin{lemma}[Exact masked attention bounds]\label{lem:reuse_attention}
\lean{OpenTheory.Attention.weight_nonneg, OpenTheory.Attention.weights_sum_one, OpenTheory.Attention.rowOutput_blocked, OpenTheory.Attention.rowOutput_mem_convexHull, OpenTheory.Attention.rowOutput_perturbation, OpenTheory.Attention.masked_output_abs_le, OpenTheory.Attention.causal_future_zero, OpenTheory.Attention.causal_row_sum_one}
\leanok
Hard-masked attention has nonnegative weights. If a row has an allowed key,
its weights sum to one and its output lies in the convex hull of its values.
With fixed weights, perturbing every value by at most $\delta$ perturbs the
output by at most $\delta$. Every actual scaled-dot-product output coordinate
is bounded in absolute value by a common bound on the corresponding value
coordinates. All-blocked rows return zero. Causal masks exclude future keys
and retain the diagonal, so every causal row is normalized.
\end{lemma}
\begin{proof}
\leanok
Reuse TorchLean's exact hard-mask coordinate formula and positive denominator;
apply Mathlib's convex-combination and norm inequalities. Identify the actual
matrix-product coordinates. This does not bound perturbations of queries or keys.
\end{proof}

\begin{definition}[Masked transformer inference]\label{def:reuse_transformer}
\lean{OpenTheory.Transformer.Config, OpenTheory.Transformer.Config.head_coverage, OpenTheory.Transformer.Config.head_dim_pos, OpenTheory.Transformer.forward, OpenTheory.Transformer.causalForward, OpenTheory.Transformer.forward_eq_upstream}
\leanok
A deterministic post-norm transformer block reuses TorchLean multi-head
attention, feedforward layers, residual addition and layer normalization.
Configuration requires positive dimensions and epsilon, and head divisibility,
ensuring positive head width and full dimension coverage. Arbitrary and causal
masks are supported. Without a mask and with the upstream epsilon, the wrapper
equals the upstream encoder. No dropout or training convergence is asserted.
\end{definition}

\begin{lemma}[Population-risk envelopes]\label{lem:reuse_risk_envelopes}
\lean{OpenTheory.Population.risk_le_envelope, OpenTheory.Population.risk_le_secondMoments, OpenTheory.Population.attentionRisk_le}
\leanok
\uses{def:reuse_data, lem:reuse_attention}
Let observations have any probability law, and assume the indicated squared
loss is almost-everywhere strongly measurable. For a $K$-Lipschitz predictor
$f$, finite input and label second moments imply integrability and
\[
\mathbb E\|f(X)-Y\|^2\leq
3\bigl(K^2\mathbb E\|X\|^2+\|f(0)\|^2+\mathbb E\|Y\|^2\bigr).
\]
For an attention row with an allowed key almost surely, values bounded by $B$
and labels bounded by $T$, the expected squared error is at most $(B+T)^2$,
even when scores and masks depend on the observation. These statements permit
arbitrary measurable observation types and normed feature spaces under their
explicit hypotheses; they are not empirical-to-population gap bounds.
\end{lemma}
\begin{proof}
\leanok
Bound the squared error pointwise using the triangle inequality and Lipschitz
or convex-combination bounds. Dominate it by an integrable envelope and apply
monotonicity and linearity of the Bochner integral.
\end{proof}

\begin{lemma}[Conditional empirical-risk oracle bridge]\label{lem:reuse_erm_bridge}
\lean{OpenTheory.Dataset.empiricalRisk_eq_foml, OpenTheory.Population.risk_eq_foml, OpenTheory.Population.bounded_loss_integrable, OpenTheory.Population.supervised_excessRisk_le}
\leanok
\uses{def:reuse_data}
Our empirical and population risks agree with FoML's risk definitions. A
measurable uniformly bounded loss is integrable under a population law.
If a trained predictor has empirical risk within $\eta$ of every predictor
in a fixed class, and every empirical/population discrepancy is at most
$\varepsilon$, its excess population risk relative to any comparator in
that class is at most $2\varepsilon+\eta$.
This is conditional: no probability bound for the deviation or optimization
guarantee for a particular neural training algorithm is asserted.
\end{lemma}
\begin{proof}
\leanok
Identify the risk normalizations and apply FoML's deterministic
approximate-ERM oracle inequality. Integrability follows by bounded domination.
\end{proof}

\begin{lemma}[Finite-class uniform deviation]\label{lem:reuse_finite_class}
\lean{OpenTheory.Population.finiteClass_deviation_le}
\leanok
\uses{def:reuse_data, lem:reuse_erm_bridge}
For a fixed finite class $\mathcal H$ of predictors whose supervised losses
are measurable and lie in $[0,1]$, and for $n\geq1$ IID observations,
\[
\Pr\!\left(\exists h\in\mathcal H:\
  |\widehat R_n(h)-R(h)|>\varepsilon\right)
\leq 2|\mathcal H|\exp(-2n\varepsilon^2)
\]
for every $\varepsilon>0$. The sample event uses the product law. The class
must be fixed independently of the sample; this theorem alone does not cover
a continuous neural parameter class.
\end{lemma}
\begin{proof}
\leanok
Identify OpenTheory's supervised risks with StatLean's definitions, then reuse
its sub-Gaussian Hoeffding bound and finite-union argument.
\end{proof}

\begin{lemma}[Ridge objective convergence]\label{lem:reuse_ridge_objective}
\lean{OpenTheory.Ridge.objective_gap_le, OpenTheory.Ridge.training_objective_gap}
\leanok
\uses{lem:reuse_ridge}
For nonnegative data weight and positive ridge regularization, exact gradient
descent satisfies the existing geometric distance bound and therefore a
geometric regularized-objective gap bound. The conversion constant is
$\frac c2\|A\|^2+\frac\lambda2$. This concerns the actual iterates and supplies
the optimization component for a later population-risk theorem.
\end{lemma}
\begin{proof}
\leanok
Expand the quadratic at the stationary minimizer, bound $\|A v\|$ by
$\|A\|\|v\|$, and compose with the checked TorchLean distance rate.
\end{proof}

\begin{lemma}[Learning by finite model selection]\label{lem:finite_learning}
\lean{OpenTheory.FiniteLearning.select, OpenTheory.FiniteLearning.select_isERM, OpenTheory.FiniteLearning.approx_select_tail, OpenTheory.FiniteLearning.select_tail}
\leanok
\uses{lem:reuse_finite_class}
For a nonempty finite collection of models fixed before sampling, an empirical
minimizer exists. For measurable losses in $[0,1]$, $n\geq1$ IID observations
and $\varepsilon>0$, the outer probability that its population risk exceeds
any candidate's by more than $2\varepsilon$ is at most
$2|\mathcal H|\exp(-2n\varepsilon^2)$. An approximate selector adds its
empirical optimization error $\eta$ to the threshold. The exact selector uses
classical choice; executable real-number comparison is not asserted.
\end{lemma}
\begin{proof}
\leanok
Use Mathlib's finite minimum theorem, then contain the excess-risk event in
StatLean's finite-class uniform-deviation event.
\end{proof}

\begin{lemma}[Continuous-class cover transfer]\label{lem:loss_cover}
\lean{OpenTheory.Covering.risk_close, OpenTheory.Covering.empirical_close, OpenTheory.Covering.uniform_deviation_tail, OpenTheory.Covering.compact_loss_cover, OpenTheory.Covering.approx_learner_tail}
\leanok
\uses{lem:reuse_finite_class, lem:finite_learning}
Let a fixed class have measurable losses in $[0,1]$ and an internal finite
uniform loss cover of size $N$ and error $\rho$. For $n\geq1$ IID observations
and $\varepsilon>0$, the outer probability of a uniform deviation exceeding
$\varepsilon+2\rho$ is at most $2N\exp(-2n\varepsilon^2)$.
An approximate empirical minimizer with error $\eta$ has the same tail bound
for excess risk above $2\varepsilon+4\rho+\eta$, against every class member.
Compact parameter classes with uniformly parameter-Lipschitz losses admit
finite covers; no dimension-dependent cardinality estimate is asserted.
\end{lemma}
\begin{proof}
\leanok
Each covered loss differs from its representative by at most $\rho$ in both
empirical and population risk. Contain the bad event in StatLean's finite-class
event. Reuse Mathlib's compact finite-ball cover for existence.
\end{proof}

\begin{lemma}[A fixed ball for ridge training]\label{lem:ridge_radius}
\lean{OpenTheory.Ridge.minimizer_norm_le, OpenTheory.Ridge.training_radius, OpenTheory.Dataset.labels_energy_le, OpenTheory.Dataset.ridge_training_radius}
\leanok
\uses{lem:reuse_ridge, lem:reuse_ridge_objective}
For $c\geq0$, $\lambda>0$, $R\geq0$ and
$c\|y\|^2\leq\lambda R^2$, every ridge minimizer has norm at most $R$.
Under the existing admissible-step condition, all exact-GD iterates from zero
have norm at most $2R$. For nonempty finite data with labels bounded by $T\geq0$,
$T^2\leq\lambda R^2$ suffices, independently of the sampled inputs.
The admissible step can still depend on the dataset.
\end{lemma}
\begin{proof}
\leanok
Compare the minimizer's objective with the objective at zero. The geometric
distance bound never exceeds the initial distance, and the triangle inequality
gives the iterate radius. Bound mean squared labels by $T^2$.
\end{proof}

\begin{lemma}[Trained ridge excess risk]\label{lem:trained_ridge}
\lean{OpenTheory.TrainedRidge.trained_ridge_tail, OpenTheory.Dataset.uniformStep, OpenTheory.Dataset.uniformFactor, OpenTheory.Dataset.uniformStep_admissible, OpenTheory.Dataset.data_q_le_uniformFactor, OpenTheory.BallCover.euclideanBall_net, OpenTheory.Ridge.minimizer_grad_eq_zero}
\leanok
\uses{lem:ridge_radius, lem:loss_cover, lem:reuse_ridge_objective}
Let $B,T,R\geq 0$, $\lambda>0$, $T^2\leq\lambda R^2$, $n\geq 1$, $\varepsilon>0$
and $\rho>0$, and let the population be supported on $\|x\|\leq B$ and
$|y|\leq T$. Run zero-initialized exact gradient descent on the ridge objective
for $k$ steps with step $\lambda/(B^2+2\lambda)^2$. On the ball of radius $2R$,
squared prediction error is at most $(2RB+T)^2$; dividing by that scale puts the
loss in $[0,1]$. Outside the bounded set the normalized loss is zero.
The outer probability that some comparator of norm at most $R$ beats the trained
predictor's normalized population risk by more than
\[
2(\varepsilon+2\rho)+\frac{(B^2+\lambda)q^k R^2+\lambda R^2}{(2RB+T)^2}
\]
is at most $2(N+1)^d\exp(-2n\varepsilon^2)$, where
$q=1-(\lambda/(B^2+2\lambda))^2$ and $N$ is the grid-cover axis count of the
Euclidean ball of radius $2R$. When the scale is zero the normalized risks
vanish and the bad event is empty. The grid cardinality is a convenient bound,
not a sharp covering number. Gradient descent minimizes the regularized
objective, so the budget pays the optimization gap and the comparator's
regularization. This is a linear ridge guarantee.
\end{lemma}
\begin{proof}
\leanok
The uniform step is admissible for every design whose rows are bounded by $B$,
and the resulting contraction is at most $q$. A minimizer is the unique
stationary point, so the distance bound and the quadratic expansion control the
regularized gap. A floor grid, radially projected onto the ball, is a finite
net; the finite-class Hoeffding bound transfers across that net. Product
samples that leave the bounded set have measure zero.
\end{proof}

\begin{lemma}[Fixed-design ridge degrees of freedom]\label{lem:ridge_risk}
\lean{OpenTheory.RidgeRisk.fixed_design_mse, OpenTheory.RidgeRisk.fixed_design_mse_of_rank, OpenTheory.RidgeRisk.degreesOfFreedom, OpenTheory.RidgeRisk.degreesOfFreedom_le_rank, OpenTheory.RidgeRisk.gd_approaches_ridgeSolution, OpenTheory.RidgeRisk.LabelNoise, OpenTheory.RidgeRisk.design_rank_le_input, OpenTheory.RidgeRisk.design_rank_le_output}
\leanok
\uses{lem:reuse_ridge_objective}
Let $m\geq 1$, $\lambda>0$, and let the design $A$ be held fixed. Labels are
$A\theta_\star+\varepsilon$, where $\varepsilon$ is white noise of scale $\sigma$:
integrable, mean zero, and with covariance $\sigma^2 I$. Exact gradient descent
on the ridge objective converges to the unique ridge minimizer. The in-sample
mean squared error of that minimizer is at most
\[
\frac{\lambda}{4}\|\theta_\star\|^2+\sigma^2\frac{\mathrm{df}(A,\lambda)}{m}.
\]
The degrees of freedom are the trace of the ridge hat operator. They are at most
$\operatorname{rank}(A)$, and that rank is at most $\min(d,m)$. If the rank is
at most $r$, the variance term may be replaced by $\sigma^2 r/m$. The bias term
is the price of regularization and does not vanish as $m$ grows. Ambient
dimension enters only through the rank, so the bound is informative when the
design rank or the degrees of freedom are much smaller than $d$. This is the
classical fixed-design bias--variance shape for ridge regression. It is not a
random-design generalization theorem and not a neural-network theorem. It is
the linear prototype of a kernel or NTK guarantee.
\end{lemma}
\begin{proof}
\leanok
A minimizer is the unique stationary point, so the gradient-descent limit equals
the ridge solution. The prediction error splits into a regularization bias and
the hat operator applied to the noise. The bias quadratic form is at most
$\lambda/(4c)$ times $\|\theta_\star\|^2$; with $c=1/m$, dividing by $m$ leaves
$(\lambda/4)\|\theta_\star\|^2$. The hat operator is self-adjoint with
eigenvalues in $[0,1]$, so its trace is at most its number of nonzero
eigenvalues, which is at most the rank of $A$. White noise has mean zero, so
the cross term integrates to zero, and the expected noise energy is $\sigma^2$
times the sum of squared eigenvalues, hence at most $\sigma^2$ times the
degrees of freedom.
\end{proof}

\begin{lemma}[Kernel form of ridge]\label{lem:kernel_ridge}
\lean{OpenTheory.KernelRidge.prediction_eq_kernelRidge, OpenTheory.KernelRidge.hat_eq_kernelSmoother, OpenTheory.KernelRidge.prediction_eq_of_kernelOp, OpenTheory.KernelRidge.kernelSmoother_eq_of_kernelOp, OpenTheory.KernelRidge.kernel_rank_eq, OpenTheory.KernelRidge.fixed_design_kernel_mse, OpenTheory.KernelRidge.kernelOp}
\leanok
\uses{lem:ridge_risk}
Let $K=cAA^*$ be the kernel operator of a fixed feature map $A$. Exact gradient
descent on the primal ridge objective has predictions $K(K+\lambda I)^{-1}y$.
The ridge hat operator is this same smoother, so the degrees of freedom in the
fixed-design bound are the trace of $K(K+\lambda I)^{-1}$. Two feature maps
with the same kernel produce the same predictions. For $c\neq 0$, the rank of
$K$ equals the rank of $A$, so the feature dimension cancels. The in-sample
mean squared error of the kernel smoother obeys the same bound as primal ridge.
The bias term still uses the norm of the primal parameter in the chosen feature
coordinates. This is the classical primal--dual identity for kernel ridge. A
frozen feature map, including Jacobian features of a linearized model, inherits
it. It does not prove that a neural network stays close to its linearization.
\end{lemma}
\begin{proof}
\leanok
The operators satisfy $(K+\lambda I)A=A(cA^*A+\lambda I)$. Inverting and
composing with $A^*$ identifies the primal predictor with the kernel smoother.
Equal kernels determine the same smoother. The range of $AA^*$ equals the range
of $A$, and a nonzero scalar does not change that range. Rewrite the
fixed-design integral through the prediction identity.
\end{proof}

\begin{lemma}[One-step feature displacement]\label{lem:feature_step}
\lean{OpenTheory.FeatureStep.preactivation_step, OpenTheory.FeatureStep.abs_preactivation_step, OpenTheory.FeatureStep.loss_hasDerivAt_entry, OpenTheory.FeatureStep.ntk_preactivation_step_le, OpenTheory.FeatureStep.mup_preactivation_step_le, OpenTheory.FeatureStep.feature_learning_preactivation_step_le, OpenTheory.FeatureStep.feature_learning_preactivation_step_ge, OpenTheory.FeatureStep.featureLearningRate}
\leanok
\uses{def:two_layer_net}
One gradient step on the half-squared loss of a two-layer network changes hidden preactivation $i$ by
\[
\Delta\mathrm{pre}_i
=
-\eta_1\,(f(x)-y)\,\sigma'(\mathrm{pre}_i)\,W_{2,i}\,\|x\|^2.
\]
If the readout coordinate is at most its initialization scale, the checked NTK rate $\eta_1=1$ moves that preactivation by $O(1/\sqrt{m})$, and the repository's checked $\mu$P rate $\eta_1=1$ moves it by $O(1/m)$. Both vanish as the width grows, so both are kernel-regime steps for this forward pass. The first-layer rate $\eta_1=m$, with a readout coordinate of size $1/m$, moves the preactivation by a width-independent amount: at most $|f(x)-y|\,L\,\|x\|^2$, and at least $c^3$ when the readout, derivative, residual and input energy are all at least $c$. This is the one-step calculation behind the Yang--Hu exponent, not the infinite-width dynamical dichotomy and not a generalization theorem. Kernel ridge applies to these features only while the displacement stays negligible.
\end{lemma}
\begin{proof}
\leanok
The partial derivative of the half-squared loss in a first-layer entry is the residual times the readout coordinate, the activation derivative and the corresponding input coordinate. Summing the resulting entrywise updates against the input produces the preactivation identity. The three rates are the absolute value of that identity under the corresponding bounds on the readout.
\end{proof}

\begin{lemma}[Tangent step and ridge perturbation]\label{lem:tangent_step}
\lean{OpenTheory.TangentStep.tangentUpdate_eq, OpenTheory.TangentStep.forward_after_sub_tangent, OpenTheory.TangentStep.abs_quadratic_remainder, OpenTheory.TangentStep.ntk_tangent_gap, OpenTheory.TangentStep.mup_tangent_gap, OpenTheory.TangentStep.feature_learning_tangent_gap, OpenTheory.TangentStep.tangentOutput, OpenTheory.PredictionGap.sq_norm_le_double, OpenTheory.PredictionGap.expected_mse_of_near_ridge}
\leanok
\uses{lem:feature_step, lem:ridge_risk}
The network output after one first-layer step equals the tangent update plus the sum of activation remainders. The tangent update is $- \eta_1 (f(x)-y) \|x\|^2 \sum_i (W_{2,i}\sigma'(\mathrm{pre}_i))^2$. If $\sigma'$ is $M$-Lipschitz, the mean value theorem bounds each activation remainder by $M$ times the squared preactivation move; the integral form would replace $M$ by $M/2$. Under the checked NTK initialization the output gap is at most $M |f(x)-y|^2 L^2 \|x\|^4 / \sqrt{m}$. Under the checked $\mu$P initialization it is at most the same quantity over $m^2$. Both vanish as the width grows. The feature-learning rate $\eta_1=m$ leaves an output gap of order $|f(x)-y|^2$, with no negative power of the width. Separately, a predictor within $\delta$ of the ridge predictor has in-sample mean squared error at most twice the ridge bound plus $2\delta^2/m$. The factor two is $2ab\le a^2+b^2$. Lemma~\ref{lem:mlp_batch_risk} discharges this $\delta$ over $T$ steps of multi-observation nonlinear batch gradient descent.
\end{lemma}
\begin{proof}
\leanok
Substitute the exact preactivation step into the first-order output formula. The difference from the updated network is a sum of activation remainders. Apply the mean value theorem on the segment from the old preactivation to the new one, and use the Lipschitz bound on $\sigma'$. Sum the resulting coordinate bounds under each learning-rate hypothesis. For the ridge comparison, expand $\|p-t\|^2$ by the triangle inequality and replace the cross term using $2ab\le a^2+b^2$. Integrate and insert the fixed-design bound.
\end{proof}

\begin{lemma}[Single-input nonlinear trajectory and centered ridge endpoint]\label{lem:mlp_single_trajectory}
\lean{OpenTheory.MLPTrajectory.preShift_step, OpenTheory.MLPTrajectory.forward_remainder_le, OpenTheory.MLPTrajectory.iterate_preShift_le, OpenTheory.MLPTrajectory.iterate_remainder_le, OpenTheory.MLPComparison.explicit_nonlinear_linearized_gap, OpenTheory.MLPEndpoint.nonlinear_endpoint_gap, OpenTheory.MLPResidual.residual_le, OpenTheory.MLPRisk.trained_endpoint_gap, OpenTheory.MLPRisk.expected_mse}
\leanok
\uses{lem:tangent_step}
For a single input $x\in\mathbb{R}^d$ with $\|x\|^2=E$, first-layer regularized gradient descent on $(f_U(x)-y)^2/2 + (\lambda/2)\|U\|_F^2$ induces a scalar recurrence on the initial-feature projection and hidden preactivations. When $|1-\eta\lambda|\le 1$, $|W_{2,i}|\le A$, $|\sigma'|\le L$, and $\sigma'$ is $M$-Lipschitz, the initial residual bound $|f_0(x)-y|\le S$ implies an a priori geometric envelope $\|U_t\|_F\le S g^t$, preactivation shift $|\Delta\mathrm{pre}_i(t)|\le S g^t\sqrt{E}$, and residual $|f_{U_t}(x)-y|\le (1+m A L\sqrt{E})S g^t$ for all $t\le T$, where $g=1+\eta\sqrt{m}A L E(1+m A L\sqrt{E})$. Accumulating the local Taylor remainder and derivative drift along the contracting linear ridge recurrence ($|1-\eta(\lambda+K)|\le q<1$) bounds the $T$-step distance between the nonlinear trained network output and the centered ridge prediction $f_0(x)+\frac{K}{K+\lambda}(y-f_0(x))$, yielding an expected squared error bound under discrete centered label noise without unproved trajectory hypotheses.
\end{lemma}
\begin{proof}
\leanok
Bound the Frobenius norm of each first-layer update by Cauchy--Schwarz and induction on $t$. Derive the preactivation shift and output residual envelopes from the parameter norm. Combine the one-step output remainder and slope drift with the scalar linear contraction recurrence to bound the $T$-step trajectory gap, add the geometric optimization residual $q^T|f_0(x)-y|$, and integrate against the discrete noise distribution.
\end{proof}

\begin{lemma}[Multi-observation nonlinear MLP trajectory and fixed-design statistical risk]\label{lem:mlp_batch_risk}
\lean{OpenTheory.MLPFixedDesign.predict_flatten, OpenTheory.MLPFixedDesign.gd_to_ridge, OpenTheory.MLPFixedDesign.ridge_expected_mse, OpenTheory.MLPBatchBounds.output_remainder_le, OpenTheory.MLPBatchBounds.design_drift_le, OpenTheory.MLPBatchBounds.adjoint_residual_gap_le, OpenTheory.MLPBatchTrajectory.batch_norm_le, OpenTheory.MLPBatchTrajectory.batch_shift_le, OpenTheory.MLPBatchTrajectory.batch_residual_le, OpenTheory.MLPBatchTrajectory.nonlinear_linear_param_gap, OpenTheory.MLPBatchTrajectory.nonlinear_linear_output_gap, OpenTheory.MLPBatchTrajectory.batch_trained_ridge_gap, OpenTheory.MLPBatchRisk.batch_trained_pointwise_error_le, OpenTheory.MLPBatchRisk.batch_trained_expected_mse, OpenTheory.MLPBatchRisk.batch_trained_expected_mse_of_rank}
\leanok
\uses{lem:mlp_single_trajectory, lem:ridge_risk}
Let $X\in\mathbb{R}^{n\times d}$ be an arbitrary $n$-observation dataset with row energies $\|X_r\|_2^2\le E$, and let $J_0=\mathrm{design}(\mathrm{net},X):\mathbb{R}^{md}\to\mathbb{R}^n$ be the multi-row initial-Jacobian feature operator around initialization $F_0=\mathrm{initialOutput}(\mathrm{net},X)$. Consider $T$ steps of exact batch gradient descent on the nonlinear first-layer regularized empirical risk $\frac{1}{2n}\|F(u)-y\|_2^2+\frac{\lambda}{2}\|u\|_2^2$ with step size $\eta>0$ satisfying $|1-\eta\lambda|\le 1$ and the ridge step condition $\eta L_{\mathrm{ridge}}^2<2\lambda$. If $\|F_0-y\|_2\le S$, then for all $t\le T$:
\begin{enumerate}
\item The parameter norm satisfies $\|u_t\|_2\le S g_n^t$, preactivation shifts satisfy $|\Delta\mathrm{pre}_i(X_r,t)|\le S g_n^t\sqrt{E}$, and batch residuals satisfy $\|F(u_t)-y\|_2\le (1+\sqrt{n}m A L\sqrt{E})S g_n^t$, where $g_n=1+\frac{\eta}{n}\sqrt{n}\sqrt{m}A L E(1+\sqrt{n}m A L\sqrt{E})$.
\item The $T$-step nonlinear batch prediction $F(u_T)$ stays within an explicit deterministic budget $\delta=\mathrm{batchGap}(T,m,n,\eta,\lambda,A,L,M,E,\|J_0\|_{\mathrm{op}},q,S)$ of the centered initial-Jacobian ridge prediction $F_0+J_0(J_0^*J_0+n\lambda I)^{-1}J_0^*(y-F_0)$.
\item When labels are generated as $y=F_0+J_0 u_\star+z$ with mean-zero variance-$\sigma^2$ label noise $z$ satisfying $\|z\|_2\le B_z$ almost surely (setting $S=\|J_0 u_\star\|_2+B_z$), the expected in-sample mean squared error of the nonlinear trained network $F(u_T)$ over the $n$ observations is at most
\[
2\left(\frac{\lambda}{4}\|u_\star\|_2^2+\sigma^2\frac{\mathrm{df}(J_0,\lambda)}{n}\right)+\frac{2\delta^2}{n}
\le
2\left(\frac{\lambda}{4}\|u_\star\|_2^2+\sigma^2\frac{\operatorname{rank}(J_0)}{n}\right)+\frac{2\delta^2}{n}.
\]
\end{enumerate}
\end{lemma}
\begin{proof}
\leanok
Flatten the $m\times d$ displacement matrix via $\mathrm{finProdFinEquiv}$ so the multi-row linearized forward pass equals $F_0+J_0 u$. Bound the operator norm of the shifted Jacobian design at every step by $\sqrt{n}\sqrt{m}A L\sqrt{E}$ and induct on $t$ to obtain the closed norm, preactivation-shift, and residual envelopes from $\|F_0-y\|_2\le S$. Decompose each parameter step difference $u_{t+1}-v_{t+1}$ into the nonexpansive linear ridge step applied to $u_t-v_t$ plus the adjoint-residual Taylor and Jacobian-drift defect $(\eta/n)(\|J_0\|_{\mathrm{op}}B+G R)$. Sum the defects over $T$ steps, convert parameter distance to output distance via $B+\|J_0\|_{\mathrm{op}}\|u_T-v_T\|_2$, and combine with geometric ridge GD convergence $\|v_T-u_\star^{\mathrm{ridge}}\|_2\le q^{T/2}S/\sqrt{n\lambda}$. Finally, apply the pointwise quadratic perturbation bound $\|p-t\|_2^2\le 2\|q-t\|_2^2+2\delta^2$ and integrate against the fixed-design ridge degrees-of-freedom theorem.
\end{proof}

\begin{lemma}[Horizon-free contractive MLP trajectory and statistical risk]\label{lem:mlp_contractive}
\lean{OpenTheory.MLPContractive.nonlinearBatchStep_norm_le, OpenTheory.MLPContractive.contractive_norm_le, OpenTheory.MLPContractive.contractive_shift_le, OpenTheory.MLPContractive.contractive_residual_le, OpenTheory.MLPContractive.accumulated_contracting_vector_defect, OpenTheory.MLPContractive.linearBatchStep_contractive, OpenTheory.MLPContractive.contractive_nonlinear_linear_param_gap, OpenTheory.MLPContractive.contractive_nonlinear_linear_output_gap, OpenTheory.MLPContractive.contractive_batch_trained_ridge_gap, OpenTheory.MLPContractive.contractive_batch_pointwise_error_le, OpenTheory.MLPContractive.contractive_batch_expected_mse, OpenTheory.MLPContractive.contractive_batch_expected_mse_of_rank}
\leanok
\uses{lem:mlp_batch_risk}
When the nonlinear contraction factor $q_{\mathrm{nl}} = |1-\eta\lambda| + \eta m\sqrt{m}A^2 L^2 E < 1$ (which holds whenever $0<\eta$, $\eta\lambda\le 1$, and $m\sqrt{m}A^2 L^2 E < \lambda$), the multi-observation nonlinear batch gradient descent trajectory satisfies horizon-free uniform envelopes for all $t\in\mathbb{N}$:
\[
\|u_t\|_2 \le \frac{c_{\mathrm{nl}} S}{1 - q_{\mathrm{nl}}} =: U_\infty,\qquad
|\Delta\mathrm{pre}_i(X_r,t)| \le U_\infty\sqrt{E} =: D_\infty,\qquad
\|F(u_t)-y\|_2 \le S + \sqrt{n}m A L D_\infty =: R_\infty.
\]
Moreover, under the ridge step condition $\eta L_{\mathrm{ridge}}^2 < 2\lambda$, the linear initial-feature step is strictly contractive with factor $q_{\mathrm{gd}} = \sqrt{q(\eta,\lambda,L_{\mathrm{ridge}})} < 1$, so the accumulated nonlinear-to-linear parameter defect is bounded uniformly over all $T\in\mathbb{N}$ by $\delta_{\mathrm{step}}/(1-q_{\mathrm{gd}})$, and the distance to the centered ridge prediction is bounded by
\[
\delta_\infty(T) = B_\infty + \|J_0\|_{\mathrm{op}}\frac{\delta_{\mathrm{step}}}{1-q_{\mathrm{gd}}} + q_{\mathrm{gd}}^T \frac{S}{\sqrt{n\lambda}}\|J_0\|_{\mathrm{op}},
\]
yielding a horizon-uniform expected in-sample MSE bound under bounded white label noise in which the optimization error decays geometrically as $T\to\infty$ without any exponential $g_n^T$ trajectory blowup.
\end{lemma}
\begin{proof}
\leanok
Combine the operator-norm bound on the shifted Jacobian with the output-change bound to establish the one-step affine norm recurrence $\|u_{t+1}\|_2 \le q_{\mathrm{nl}}\|u_t\|_2 + c_{\mathrm{nl}}S$. Induct on $t$ using $q_{\mathrm{nl}}<1$ to obtain the geometric fixed-point bound $U_\infty$, and deduce the time-uniform shift and residual envelopes $D_\infty$ and $R_\infty$. Use strict contraction of the strongly convex ridge GD step (`Optim.GD.step_norm_sq_le` and `Optim.GD.q_lt_one_of_mul_sq_lt`) to sum the one-step defects as a convergent geometric series bounded by $\delta_{\mathrm{step}}/(1-q_{\mathrm{gd}})$. Combine with geometric convergence to the ridge solution and integrate via `PredictionGap.sq_norm_le_double` and `MLPFixedDesign.ridge_expected_mse`.
\end{proof}

\begin{lemma}[Feature-learning advantage over frozen initial kernels]\label{lem:feature_learning_advantage}
\lean{OpenTheory.FeatureLearningAdvantage.feature_learning_preactivation_shift_exact, OpenTheory.FeatureLearningAdvantage.standard_lr_preactivation_shift_suppressed, OpenTheory.FeatureLearningAdvantage.feature_learning_vs_standard_shift_ratio, OpenTheory.FeatureLearningAdvantage.feature_learning_preactivation_sq_strict_growth, OpenTheory.FeatureLearningAdvantage.kernel_strict_growth_of_slope_growth, OpenTheory.FeatureLearningAdvantage.ridgeWeights_objective_eq_minRidgeObj, OpenTheory.FeatureLearningAdvantage.endpoint_expected_mse_eq_ridgeExpectedMSE, OpenTheory.FeatureLearningAdvantage.minRidgeObj_strict_anti, OpenTheory.FeatureLearningAdvantage.biasSq_sub_eq, OpenTheory.FeatureLearningAdvantage.biasSq_strict_anti, OpenTheory.FeatureLearningAdvantage.biasSq_reduction_lower_bound, OpenTheory.FeatureLearningAdvantage.ridge_mse_strict_improvement_of_kernel_growth, OpenTheory.FeatureLearningAdvantage.ridge_mse_strict_improvement_of_weighted_snr, OpenTheory.FeatureLearningAdvantage.frobeniusSq_residual_eq_one_sub_alignment, OpenTheory.FeatureLearningAdvantage.residual_projection_strict_decrease_of_alignment_growth}
\leanok
\uses{lem:feature_step, lem:mlp_single_trajectory, lem:alignment_certificate}
Under the feature-learning rate $\eta_1 = m$ and readout scale $W_{2,i} = s_i/m$ ($s_i^2=1$) at zero initial prediction $f_0(x)=0$ with local activation slope $\alpha>0$, one gradient step on target $y\ne 0$ shifts hidden preactivation $i$ by the exact width-independent quantity $\Delta\mathrm{pre}_i = y\alpha s_i\|x\|_2^2$ ($m$ times larger than the $O(1/m)$ shift under standard $\mu$P unit learning rate), strictly increasing squared preactivations and the post-step tangent kernel energy from $K_0$ to $K_1 > K_0$ when active slopes grow. For any centered ridge problem with regularization $\lambda>0$ and target offset $\Delta y\ne 0$, increasing the feature kernel energy from $K_0$ to $K_1 > K_0$ strictly decreases both the minimum regularized ridge objective $\frac{1}{2}\frac{\lambda}{K+\lambda}\Delta y^2$ and the squared regularization bias $\bigl(\frac{\lambda}{K+\lambda}\bigr)^2\Delta y^2$ by the exact gap
\[
\mathrm{biasSq}(K_0,\lambda,\Delta y) - \mathrm{biasSq}(K_1,\lambda,\Delta y)
= \frac{\lambda^2 \Delta y^2 (K_1 - K_0)(K_1 + K_0 + 2\lambda)}{(K_0+\lambda)^2(K_1+\lambda)^2} > 0,
\]
and strictly reduces the total expected ridge mean squared error $\bigl(\frac{\lambda}{K+\lambda}\bigr)^2\Delta y^2 + \bigl(\frac{K}{K+\lambda}\bigr)^2\sigma^2$ whenever $(K_1(K_0+\lambda)+K_0(K_1+\lambda))\sigma^2 < \lambda(K_1+K_0+2\lambda)\Delta y^2$ (in particular whenever $K_1\sigma^2 \le \lambda\Delta y^2$). Equivalently, at fixed total Frobenius feature energy, strictly increasing subspace alignment $a(H_0,P) < a(H_1,P)$ strictly decreases residual projection error $\lVert H_1 - H_1 P\rVert_F^2 < \lVert H_0 - H_0 P\rVert_F^2$.
\end{lemma}
\begin{proof}
\leanok
Specialize `FeatureStep.preactivation_step` to $f_0(x)=0$, $W_{2,i}=s_i/m$, and $\eta_1=m$ to cancel $m$ and obtain the exact shift $y\alpha s_i\|x\|_2^2$. Expand $(p_i + \Delta\mathrm{pre}_i)^2 - p_i^2$ and sum over neurons to establish strict kernel growth. Evaluate the centered ridge objective and finite-noise expected MSE at `MLPEndpoint.ridgeWeights`, put the difference of rational functions over the common denominator $(K_0+\lambda)^2(K_1+\lambda)^2$, and apply positivity under the signal-to-noise threshold. Combine with the Frobenius Pythagorean identity from Lemma~\ref{lem:alignment_le_one} for the subspace residual reduction.
\end{proof}

\begin{lemma}[Multi-head attention ablation fidelity, residual propagation, and risk transfer]\label{lem:attention_ablation}
\lean{OpenTheory.AttentionAblation.headContribution_norm_le, OpenTheory.AttentionAblation.ablation_diff_eq, OpenTheory.AttentionAblation.subset_ablation_norm_le, OpenTheory.AttentionAblation.single_head_ablation_norm_le, OpenTheory.AttentionAblation.multiHead_value_perturbation, OpenTheory.AttentionAblation.residual_block_ablation_bound, OpenTheory.AttentionAblation.multilayer_residual_perturbation, OpenTheory.AttentionAblation.ablation_pointwise_sq_error_le, OpenTheory.AttentionAblation.ablation_population_mse_le, OpenTheory.AttentionAblation.single_head_ablation_population_mse_le}
\leanok
\uses{lem:tangent_step}
Consider an $H$-head hard-masked attention layer $\mathrm{MHA}(x) = \sum_{h=1}^H W_{O,h}\,\mathrm{rowOutput}(s_h, m_h, v_h)$ and the head-ablation intervention $\mathrm{MHA}_{\mathrm{keep}}(x) = \sum_{h\in\mathrm{keep}} W_{O,h}\,\mathrm{rowOutput}(s_h, m_h, v_h)$ that zeros out heads in $\mathrm{keep}^c$ without rescaling remaining heads:
\begin{enumerate}
\item \textbf{Head-ablation fidelity:} If every ablated head $h\in\mathrm{keep}^c$ has at least one unmasked key and value norm bounded by $\|v_{h,i}\|\le B_{V,h}$, then
\[
\bigl\|\mathrm{MHA}(x) - \mathrm{MHA}_{\mathrm{keep}}(x)\bigr\| \le \sum_{h\in\mathrm{keep}^c} \|W_{O,h}\|_{\mathrm{op}}\,B_{V,h},
\]
and zeroing a single head $h_0$ changes the output by at most $\|W_{O,h_0}\|_{\mathrm{op}}\,B_{V,h_0}$.
\item \textbf{Downstream residual & multilayer propagation:} Passing the post-attention residual $x + \mathrm{MHA}(x)$ through an $L_{\mathrm{ff}}$-Lipschitz downstream block $G$ bounds the output change by $L_{\mathrm{ff}}\sum_{h\in\mathrm{keep}^c}\|W_{O,h}\|_{\mathrm{op}}B_{V,h}$, and across an $L$-layer residual stack with branch Lipschitz constants $c_t$ and layerwise perturbations $\varepsilon_t$, the end-to-end perturbation is bounded by the recurrence $E_{t+1} = (1+c_t)E_t + \varepsilon_t$.
\item \textbf{Statistical population risk transfer:} Under a data population $P$ where the mask and value bounds hold almost surely on $\mathrm{keep}^c$, the population squared-error risk of the ablated model is bounded by
\[
R_P(\mathrm{MHA}_{\mathrm{keep}}) \le 2\,R_P(\mathrm{MHA}) + 2\left(\sum_{h\in\mathrm{keep}^c} \|W_{O,h}\|_{\mathrm{op}}\,B_{V,h}\right)^2.
\]
\end{enumerate}
\end{lemma}
\begin{proof}
\leanok
Split the full sum over $\mathrm{Fin}\,H$ into $\mathrm{keep}$ and $\mathrm{keep}^c$ via `Finset.sum_add_sum_compl`. Bound each ablated head's row output by $B_{V,h}$ using `Attention.rowOutput_norm_le` and apply the continuous linear operator norm bound $\|W_{O,h}(u)\| \le \|W_{O,h}\|_{\mathrm{op}}\|u\|$. Propagate through $G$ by Lipschitz continuity and through $L$ residual stages by induction on $L$. For the population risk bound, apply `PredictionGap.sq_norm_le_double` pointwise a.e. and integrate against $P.\mathrm{law}$.
\end{proof}

\begin{lemma}[Local-sharpness adaptive descent, objective coercivity, and discrete PL convergence]\label{lem:local_sharpness_descent}
\lean{OpenTheory.LocalSharpnessDescent.local_sharpness_descent, OpenTheory.LocalSharpnessDescent.adaptive_step_progress_eq, OpenTheory.LocalSharpnessDescent.stepProgress_sub_eq_sq, OpenTheory.LocalSharpnessDescent.adaptive_step_maximizes_progress, OpenTheory.LocalSharpnessDescent.adaptive_step_optimal_drop, OpenTheory.LocalSharpnessDescent.adaptive_vs_global_sharpness_speedup, OpenTheory.LocalSharpnessDescent.adaptive_step_strictly_beats_conservative, OpenTheory.LocalSharpnessDescent.regularized_loss_coercive_residual_le, OpenTheory.LocalSharpnessDescent.regularized_loss_coercive_norm_le, OpenTheory.LocalSharpnessDescent.regularized_loss_coercive_shift_le, OpenTheory.LocalSharpnessDescent.descent_trajectory_horizon_free_bounds, OpenTheory.LocalSharpnessDescent.pl_adaptive_one_step_contraction, OpenTheory.LocalSharpnessDescent.pl_trajectory_geometric_convergence}
\leanok
\uses{lem:mlp_batch_risk, lem:risk_exponential}
Let $L(u)$ satisfy the local quadratic sharpness bound $L(u - \eta g) \le L(u) - \eta\|g\|_2^2 + \frac{S_{\mathrm{loc}}}{2}\eta^2\|g\|_2^2$ with local sharpness $0 < S_{\mathrm{loc}} < S_{\mathrm{global}}$:
\begin{enumerate}
\item \textbf{Optimal adaptive step size:} The guaranteed one-step loss reduction $\mathrm{stepProgress}(\eta, S_{\mathrm{loc}}, \|g\|_2^2) = \eta(1 - \eta S_{\mathrm{loc}}/2)\|g\|_2^2$ satisfies the exact completion-of-squares identity
\[
\frac{\|g\|_2^2}{2 S_{\mathrm{loc}}} - \mathrm{stepProgress}(\eta, S_{\mathrm{loc}}, \|g\|_2^2) = \frac{S_{\mathrm{loc}}}{2}\Bigl(\eta - \frac{1}{S_{\mathrm{loc}}}\Bigr)^2 \|g\|_2^2 \ge 0,
\]
so $\eta^\star = 1/S_{\mathrm{loc}}$ uniquely maximizes guaranteed progress and achieves an exact $S_{\mathrm{global}}/S_{\mathrm{loc}} > 1$ speedup over the conservative global step $1/S_{\mathrm{global}}$.
\item \textbf{Coercivity without contraction ($q_{\mathrm{nl}} < 1$ not required):} Along any monotone loss-descent trajectory $\mathcal{L}(u_t) \le \mathcal{L}(0)$ on the nonlinear regularized empirical risk $\mathcal{L}(u) = \frac{1}{2n}\|F(u)-y\|_2^2 + \frac{\lambda}{2}\|u\|_2^2$ with $\|F(0)-y\|_2\le S$, every iterate $t\in\mathbb{N}$ satisfies the horizon-free bounds
\[
\|F(u_t)-y\|_2 \le S,\qquad \|u_t\|_2 \le \frac{S}{\sqrt{n\lambda}},\qquad |\Delta\mathrm{pre}_i(X_r, t)| \le \frac{S\sqrt{E}}{\sqrt{n\lambda}}.
\]
\item \textbf{Discrete Polyak--{\L}ojasiewicz convergence:} Under the local PL condition $\|g_t\|_2^2 \ge 2\mu(L(u_t)-F)$ and $0 < \mu \le S_{\mathrm{loc}}(t) \le S_{\max}$, adaptive steps $\eta_t = 1/S_{\mathrm{loc}}(t)$ satisfy $L(u_T) - F \le (1 - \mu/S_{\max})^T (L(u_0) - F)$ for all $T\in\mathbb{N}$.
\end{enumerate}
\end{lemma}
\begin{proof}
\leanok
Complete the square in $\eta$ to prove optimality of $\eta^\star = 1/S_{\mathrm{loc}}$. For coercivity, drop each nonnegative term in $\frac{1}{2n}\|F(u_t)-y\|_2^2 + \frac{\lambda}{2}\|u_t\|_2^2 \le \frac{S^2}{2n}$ to bound $\|F(u_t)-y\|_2$ and $\|u_t\|_2$ directly, and apply Cauchy--Schwarz for the preactivation shift. For PL convergence, substitute $\|g_t\|_2^2 \ge 2\mu(L(u_t)-F)$ into the optimal drop $L(u_{t+1}) \le L(u_t) - \|g_t\|_2^2/(2 S_{\mathrm{loc}}(t))$ and induct on $T$.
\end{proof}

\begin{lemma}[Trajectory-localized out-of-sample generalization bound]\label{lem:trajectory_generalization}
\lean{OpenTheory.TrajectoryGeneralization.forward_param_diff_le, OpenTheory.TrajectoryGeneralization.paramLoss_lipschitz, OpenTheory.TrajectoryGeneralization.clippedSqLoss_lipschitz, OpenTheory.TrajectoryGeneralization.clippedSqLoss_param_lipschitz, OpenTheory.TrajectoryGeneralization.trajectoryBall_loss_cover, OpenTheory.TrajectoryGeneralization.trajectory_uniform_deviation_tail, OpenTheory.TrajectoryGeneralization.contractive_trajectory_in_ball, OpenTheory.TrajectoryGeneralization.contractive_trajectory_generalization_tail, OpenTheory.TrajectoryGeneralization.approx_erm_trajectory_excess_risk_tail}
\leanok
\uses{lem:mlp_contractive, lem:forward_lipschitz}
For a two-layer network with $|W_{2,i}|\le A$, $|\sigma'|\le L$, inputs with $\|x\|_2^2\le E$, and an $L_\ell$-Lipschitz $[0,1]$-bounded loss $\ell$ (such as the clipped squared loss $\min(1, ((p-y)/B)^2)$ with $L_\ell = 2/B$), the parameter-to-loss map $u\mapsto \ell(f_{u}(x), y)$ is $L_{\mathrm{param}}$-Lipschitz on $\mathbb{R}^{md}$ with $L_{\mathrm{param}} = L_\ell m A L\sqrt{E}$. Consequently, for any trajectory radius $R\ge 0$ (such as the contractive radius $U_\infty$ or coercive radius $S/\sqrt{n\lambda}$) and grid resolution $\delta>0$, the parameter ball $B_{md}(R) = \{u\in\mathbb{R}^{md} : \|u\|_2\le R\}$ admits an explicit finite loss cover of cardinality $N(R,\delta) \le (\mathrm{coverAxis}(R,\delta,md)+1)^{md}$ such that for every $n\ge 1$ and $\varepsilon>0$:
\[
\mathbb{P}_{s\sim P^n}\!\left(\sup_{u\in B_{md}(R)} \bigl|\widehat{R}_s(u) - R_P(u)\bigr| > \varepsilon + 2 L_{\mathrm{param}}\delta\right)
\le 2\,\bigl(\mathrm{coverAxis}(R,\delta,md)+1\bigr)^{md}\exp(-2n\varepsilon^2).
\]
In particular, under the contractive training condition $q_{\mathrm{nl}}<1$, every gradient descent iterate $u_T$ ($T\in\mathbb{N}$) lies in $B_{md}(U_\infty)$ and satisfies this out-of-sample generalization bound uniformly in $T$.
\end{lemma}
\begin{proof}
\leanok
Write $f_u(x) - f_v(x)$ as the forward change from base network $\mathrm{atDisplacement}(\mathrm{net}, v)$ under displacement $u-v$, bound preactivation shifts by $\|u-v\|_2\sqrt{E}$, and apply `MLPResidual.forward_change_le` to obtain $|f_u(x)-f_v(x)| \le m A L\sqrt{E}\|u-v\|_2$. Combine with `BallCover.euclideanBall_net` to construct the finite loss cover of $B_{md}(R)$ and apply `Covering.uniform_deviation_tail` and `Covering.approx_learner_tail`.
\end{proof}

\begin{lemma}[Softmax score sensitivity, joint Q/K/V stability, and logit-margin concentration]\label{lem:attention_logit_stability}
\lean{OpenTheory.AttentionLogitStability.rowOutput_sub_eq_centered_sum, OpenTheory.AttentionLogitStability.rowOutput_score_perturbation_l1, OpenTheory.AttentionLogitStability.rowOutput_score_perturbation_centered, OpenTheory.AttentionLogitStability.rowOutput_joint_qkv_perturbation, OpenTheory.AttentionLogitStability.multiHead_joint_qkv_perturbation, OpenTheory.AttentionLogitStability.tail_mass_eq_of_focused, OpenTheory.AttentionLogitStability.focused_rowOutput_approx_dominant, OpenTheory.AttentionLogitStability.focused_weight_l1_bound, OpenTheory.AttentionLogitStability.focused_score_perturbation_stability, OpenTheory.AttentionLogitStability.logit_margin_implies_focused, OpenTheory.AttentionLogitStability.logit_margin_output_stability}
\leanok
\uses{lem:attention_ablation}
Let $p(s, m) = \mathrm{rowWeights}(s, m)$ denote the hard-masked softmax attention distribution on $\mathrm{Fin}\,n$ when at least one key is unmasked:
\begin{enumerate}
\item \textbf{Centered value-dispersion & joint Q/K/V perturbation:} For any reference vector $\bar{v}\in E$ with $\|v_i - \bar{v}\|\le \Delta_V$ on unmasked keys and value shift $\|v_{1,i} - v_{2,i}\|\le \delta_V$, perturbing scores $s_1\to s_2$ and values $v_1\to v_2$ satisfies
\[
\bigl\|\mathrm{rowOutput}(s_1, m, v_1) - \mathrm{rowOutput}(s_2, m, v_2)\bigr\| \le \Delta_V\,\|p(s_1,m) - p(s_2,m)\|_1 + \delta_V,
\]
and across $H$ heads, $\|\mathrm{MHA}(s_1, v_1) - \mathrm{MHA}(s_2, v_2)\| \le \sum_{h=1}^H \|W_{O,h}\|_{\mathrm{op}}(\Delta_{V,h}\|p_1^{(h)} - p_2^{(h)}\|_1 + \delta_{V,h})$.
\item \textbf{Focused / low-entropy stability:} If $p(s,m)_{i_\star} \ge 1 - \varepsilon$ and $\|v_j - v_{i_\star}\|\le D_V$ on unmasked keys, then $\|\mathrm{rowOutput}(s,m,v) - v_{i_\star}\| \le \varepsilon D_V$. If two score vectors focus on $i_\star$ with tail defects $\varepsilon_1, \varepsilon_2$, then $\|p(s_1,m) - p(s_2,m)\|_1 \le 2(\varepsilon_1 + \varepsilon_2)$ and $\|\mathrm{rowOutput}(s_1,m,v) - \mathrm{rowOutput}(s_2,m,v)\| \le (\varepsilon_1 + \varepsilon_2)D_V$.
\item \textbf{Logit-margin exponential concentration:} If an unmasked key $i_\star$ has logit margin $s_j \le s_{i_\star} - \gamma$ over all other unmasked keys $j\ne i_\star$, then $p(s,m)_j \le e^{-\gamma} p(s,m)_{i_\star}$, the tail mass satisfies $1 - p(s,m)_{i_\star} \le (n-1)e^{-\gamma}$, and $\|\mathrm{rowOutput}(s,m,v) - v_{i_\star}\| \le (n-1)e^{-\gamma} D_V$.
\end{enumerate}
\end{lemma}
\begin{proof}
\leanok
Since $\sum_i p(s_1,m)_i = \sum_i p(s_2,m)_i = 1$, subtracting $\sum_i (p_1(i)-p_2(i))\bar{v} = 0$ rewrites the output difference in terms of centered values $v_i - \bar{v}$, yielding the $\Delta_V\|p_1-p_2\|_1$ bound by the triangle inequality. For a dominant key $i_\star$, rewrite $\mathrm{rowOutput}(s,m,v) - v_{i_\star} = \sum_{j\ne i_\star} p_j(v_j - v_{i_\star})$ and use $\sum_{j\ne i_\star} p_j = 1 - p_{i_\star} \le \varepsilon$. Under a logit margin $\gamma$, `Attention.weight_formula` gives $\exp(s_j) \le e^{-\gamma}\exp(s_{i_\star})$ on unmasked keys $j\ne i_\star$, and summing over the $n-1$ non-dominant keys bounds the tail mass by $(n-1)e^{-\gamma}$.
\end{proof}

\begin{lemma}[First-principles MLP local sharpness and unconditional coercive trajectory]\label{lem:mlp_derived_sharpness}
\lean{OpenTheory.MLPLocalSharpness.nonlinearBatchStep_eq_sub_batchGrad, OpenTheory.MLPLocalSharpness.mlpLocalSharpness_eq_of_pos, OpenTheory.MLPLocalSharpness.stepOutputChange_norm_le, OpenTheory.MLPLocalSharpness.stepRemainder_norm_le, OpenTheory.MLPLocalSharpness.mlp_objective_quadratic_upper_bound, OpenTheory.MLPLocalSharpness.mlp_local_sharpness_step_descent, OpenTheory.MLPLocalSharpness.mlp_adaptive_step_optimal_drop, OpenTheory.MLPLocalSharpness.mlp_safe_step_objective_descent, OpenTheory.MLPLocalSharpness.mlp_safe_step_coercive_trajectory, OpenTheory.MLPLocalSharpness.mlp_adaptive_step_cumulative_grad_sq_le, OpenTheory.MLPLocalSharpness.mlp_adaptive_step_min_grad_norm}
\leanok
\uses{lem:local_sharpness_descent, lem:mlp_batch_risk}
For a two-layer network with $|W_{2,i}|\le A$, $|\sigma'|\le L$, $M$-Lipschitz $\sigma'$, and row energies $\|X_r\|_2^2\le E$, the nonlinear regularized empirical risk $\mathcal{L}(u) = \frac{1}{2n}\|F(u)-y\|_2^2 + \frac{\lambda}{2}\|u\|_2^2$ satisfies at any parameter $u$ with $\|F(u)-y\|_2\le R$ and any displacement $\Delta\in\mathbb{R}^{md}$ the exact second-order Taylor upper bound
\[
\mathcal{L}(u+\Delta) \le \mathcal{L}(u) + \langle \nabla\mathcal{L}(u), \Delta\rangle + \frac{S_{\mathrm{MLP}}(R)}{2}\|\Delta\|_2^2,
\qquad
S_{\mathrm{MLP}}(R) = \lambda + m^2 A^2 L^2 E + \frac{2R}{\sqrt{n}} M m A E.
\]
Consequently, when $\|F(0)-y\|_2\le S$ and the step size satisfies $0\le \eta \le 2/S_{\mathrm{MLP}}(S)$ (in particular at the adaptive step $\eta^\star = 1/S_{\mathrm{MLP}}(S)$), the actual nonlinear gradient descent trajectory $u_{t+1} = u_t - \eta\nabla\mathcal{L}(u_t)$ satisfies for all $t\in\mathbb{N}$ (without requiring $q_{\mathrm{nl}}<1$):
\[
\mathcal{L}(u_{t+1}) \le \mathcal{L}(u_t) \le \mathcal{L}(0),\qquad
\|F(u_t)-y\|_2 \le S,\qquad
\|u_t\|_2 \le \frac{S}{\sqrt{n\lambda}},\qquad
\min_{0\le t < T} \|\nabla\mathcal{L}(u_t)\|_2^2 \le \frac{S_{\mathrm{MLP}}(S)\,S^2}{T n}.
\]
\end{lemma}
\begin{proof}
\leanok
Decompose $F(u+\Delta) - y = r_u + (J_u\Delta + \mathrm{rem})$ around the shifted network $\mathrm{atDisplacement}(\mathrm{net}, u)$. Bound $\|J_u\Delta + \mathrm{rem}\|_2 \le \sqrt{n} m A L\sqrt{E}\|\Delta\|_2$ via `MLPBatchBounds.batchOutput_change_le` and $\|\mathrm{rem}\|_2 \le \sqrt{n} M m A E\|\Delta\|_2^2$ via `MLPBatchBounds.output_remainder_le`. Expand $\|r_u + (J_u\Delta + \mathrm{rem})\|_2^2$ and $\|u+\Delta\|_2^2$ using `norm_add_sq_real`, identify the linear term $\langle \nabla\mathcal{L}(u), \Delta\rangle$ via adjointness of $J_u$, and bound $\langle r_u, \mathrm{rem}\rangle \le R\|\mathrm{rem}\|_2$ by Cauchy--Schwarz. Induct on $t$ using Lemma~\ref{lem:local_sharpness_descent}: at each step $\mathcal{L}(u_t)\le \mathcal{L}(0)$ implies $\|r_{u_t}\|_2\le S$, which keeps the local sharpness bounded by $S_{\mathrm{MLP}}(S)$ for step $t+1$. Telescope the one-step drops to bound the minimum squared gradient norm over $T$ steps.
\end{proof}

\begin{lemma}[Diagonal preconditioned and SignGD/Adam-style optimization speedup]\label{lem:preconditioned_descent}
\lean{OpenTheory.PreconditionedDescent.anisotropic_sharpness_descent, OpenTheory.PreconditionedDescent.optimal_diagonal_progress_eq, OpenTheory.PreconditionedDescent.optimal_diagonal_descent_drop, OpenTheory.PreconditionedDescent.coordProgress_sub_eq_sum_sq, OpenTheory.PreconditionedDescent.optimal_diagonal_maximizes_progress, OpenTheory.PreconditionedDescent.diagonal_vs_scalar_gap_eq, OpenTheory.PreconditionedDescent.diagonal_dominates_scalar_sgd, OpenTheory.PreconditionedDescent.diagonal_strictly_beats_scalar_sgd, OpenTheory.PreconditionedDescent.diagonal_dominates_any_scalar_step, OpenTheory.PreconditionedDescent.signGD_descent, OpenTheory.PreconditionedDescent.signGD_optimal_lr_drop, OpenTheory.PreconditionedDescent.signGD_strictly_beats_scalar_sgd, OpenTheory.PreconditionedDescent.preconditioned_pl_geometric_convergence}
\leanok
\uses{lem:local_sharpness_descent}
Under the coordinate-wise anisotropic sharpness bound $L(u - \eta\odot g) \le L(u) - \sum_{i=1}^p \eta_i g_i^2 + \frac{1}{2}\sum_{i=1}^p S_i \eta_i^2 g_i^2$ with $0 < S_i \le S_{\max}$:
\begin{enumerate}
\item \textbf{Optimal diagonal preconditioning vs.\ scalar SGD:} The diagonal preconditioner $\eta_i^\star = 1/S_i$ uniquely maximizes coordinate progress $\mathrm{coordProgress}(\eta, S, g) = \sum_i \eta_i(1 - \eta_i S_i/2)g_i^2$, achieving $\sum_i g_i^2/(2S_i)$ and beating scalar SGD ($\eta_{\mathrm{scalar}} = 1/S_{\max}$, progress $\|g\|_2^2/(2S_{\max})$) by the exact gap
\[
\sum_{i=1}^p \frac{g_i^2}{2S_i} - \frac{\|g\|_2^2}{2S_{\max}} = \sum_{i=1}^p \frac{S_{\max} - S_i}{2 S_i S_{\max}}\,g_i^2 \ge 0,
\]
which is strictly positive whenever any active gradient coordinate $g_{i_0}\ne 0$ has $S_{i_0} < S_{\max}$.
\item \textbf{SignGD / memoryless Adam $\ell_1$ descent:} Stepping by $u_i - \eta\,\mathrm{sign}(g_i)$ at optimal learning rate $\eta_{\mathrm{sign}} = \|g\|_1 / \|S\|_1$ guarantees $L(u_{\mathrm{next}}) \le L(u) - \|g\|_1^2 / (2\|S\|_1)$, strictly beating scalar SGD whenever $\|g\|_1^2 / \|S\|_1 > \|g\|_2^2 / S_{\max}$ (dense gradients under anisotropic curvature).
\item \textbf{Preconditioned PL convergence:} Under $\sum_i g_{t,i}^2 / S_i \ge 2\mu_P(L(u_t)-F)$ with $0 < \mu_P \le 1$, diagonal preconditioned steps satisfy $L(u_T) - F \le (1-\mu_P)^T(L(u_0)-F)$.
\end{enumerate}
\end{lemma}
\begin{proof}
\leanok
Complete the square coordinate by coordinate to establish $\sum_i g_i^2/(2S_i) - \mathrm{coordProgress}(\eta,S,g) = \sum_i \frac{S_i}{2}(\eta_i - 1/S_i)^2 g_i^2 \ge 0$. Combine fractions over $2 S_i S_{\max}$ for the scalar SGD comparison and apply `Finset.sum_pos'` when an active coordinate has $S_{i_0} < S_{\max}$. For SignGD, use $\mathrm{sign}(x)x = |x|$ and $\mathrm{sign}(x)^2 \le 1$ to obtain the quadratic bound in $\eta$ and substitute $\eta_{\mathrm{sign}} = \|g\|_1/\|S\|_1$.
\end{proof}

\begin{lemma}[Nonlinear global PL convergence to zero training loss and $\mup$ acceleration]\label{lem:joint_pl_convergence}
\lean{OpenTheory.JointPLConvergence.unregMSE_eq, OpenTheory.JointPLConvergence.mlp_unreg_step_descent, OpenTheory.JointPLConvergence.jacobian_coercive_implies_pl, OpenTheory.JointPLConvergence.joint_two_layer_pl_inequality, OpenTheory.JointPLConvergence.joint_jacobian_coercive_implies_pl, OpenTheory.JointPLConvergence.mlp_unreg_adaptive_pl_step, OpenTheory.JointPLConvergence.mlp_global_geometric_convergence_to_zero_loss, OpenTheory.JointPLConvergence.mlp_global_residual_sq_geometric_decay, OpenTheory.JointPLConvergence.feature_learning_contraction_rate_strict_improvement, OpenTheory.JointPLConvergence.joint_vs_single_layer_contraction_strict_improvement, OpenTheory.JointPLConvergence.spectral_amplification_strict_pl_acceleration}
\leanok
\uses{lem:mlp_derived_sharpness, lem:spectral_bias_mup}
For the unregularized empirical MSE $\mathcal{L}_0(u) = \frac{1}{2n}\|F(u)-y\|_2^2$ of a nonlinear two-layer MLP:
\begin{enumerate}
\item \textbf{Single-layer and joint $(W_1, W_2)$ Polyak--{\L}ojasiewicz inequalities:} Whenever the local Jacobian satisfies $\|J_u^* r_u\|_2 \ge \sigma_{\min}\|r_u\|_2$, the gradient satisfies $\|\nabla\mathcal{L}_0(u)\|_2^2 \ge \frac{2\sigma_{\min}^2}{n}\mathcal{L}_0(u)$. In joint two-layer training where hidden and readout Jacobians satisfy $\|J_{W_1}^* r_u\|_2 \ge \sigma_1\|r_u\|_2$ and $\|\Phi_u^* r_u\|_2 \ge \sigma_2\|r_u\|_2$, the joint gradient satisfies $\|\nabla_{W_1}\mathcal{L}_0\|_2^2 + \|\nabla_{W_2}\mathcal{L}_0\|_2^2 \ge \frac{2(\sigma_1^2 + \sigma_2^2)}{n}\mathcal{L}_0(u)$.
\item \textbf{Global geometric convergence to zero training loss:} If $\|F(0)-y\|_2\le S$ and $\|J_{u_t}^* r_{u_t}\|_2 \ge \sigma_{\min}\|r_{u_t}\|_2$ along the adaptive-step trajectory $u_{t+1} = u_t - \frac{1}{S_{\mathrm{MLP}}(0,S)}\nabla\mathcal{L}_0(u_t)$, then for all $T\in\mathbb{N}$,
\[
\mathcal{L}_0(u_T) \le \Bigl(1 - \frac{\sigma_{\min}^2}{n\,S_{\mathrm{MLP}}(0,S)}\Bigr)^T \mathcal{L}_0(0),\qquad \|F(u_T)-y\|_2 \le S.
\]
\item \textbf{$\mup$ feature-learning acceleration:} Amplifying the effective singular value from $\sigma_0$ to $\sigma_{\mathrm{feat}} > \sigma_0$ strictly reduces the per-step contraction factor $1 - \frac{\sigma_{\mathrm{feat}}^2}{n S_{\mathrm{safe}}} < 1 - \frac{\sigma_0^2}{n S_{\mathrm{safe}}}$ and yields strictly smaller training error at every step $T\ge 1$.
\end{enumerate}
\end{lemma}
\begin{proof}
\leanok
Square $\|\nabla\mathcal{L}_0(u)\|_2 = \frac{1}{n}\|J_u^* r_u\|_2 \ge \frac{\sigma_{\min}}{n}\|r_u\|_2$ and substitute $\|r_u\|_2^2 = 2n\mathcal{L}_0(u)$. Induct on $T$: at step $T$, $\|r_{u_T}\|_2\le S$ activates \cref{lem:mlp_derived_sharpness} at $\lambda=0$ with sharpness $S_{\mathrm{MLP}}(0,S)$, giving $\mathcal{L}_0(u_{T+1}) \le (1 - \frac{\sigma_{\min}^2}{n S_{\mathrm{MLP}}(0,S)})\mathcal{L}_0(u_T) \le \mathcal{L}_0(0)$, which preserves $\|r_{u_{T+1}}\|_2\le S$.
\end{proof}

\begin{lemma}[In-context associative retrieval dynamics and step complexity]\label{lem:in_context_retrieval}
\lean{OpenTheory.InContextRetrieval.logitMargin_after_step, OpenTheory.InContextRetrieval.logitMargin_trajectory_ge, OpenTheory.InContextRetrieval.bilinearScore_step_eq, OpenTheory.InContextRetrieval.queryStep_score_separation, OpenTheory.InContextRetrieval.bilinearQueryStep_margin_growth, OpenTheory.InContextRetrieval.icl_attention_weight_convergence_real_sub, OpenTheory.InContextRetrieval.icl_retrieval_error_exponential_decay, OpenTheory.InContextRetrieval.icl_step_complexity_certificate, OpenTheory.InContextRetrieval.icl_population_risk_bound}
\leanok
\uses{lem:attention_logit_stability, lem:attention_ablation}
For a hard-masked attention head with initial logit margin $\gamma_0$ at target context token $i_\star$ ($s_{0,j} \le s_{0,i_\star} - \gamma_0$ for unmasked $j\ne i_\star$) updated by score steps $s_{t+1} = s_t + \eta g_t$ with key separation $g_{t,j} \le g_{t,i_\star} - \Delta_g$ ($\eta\ge 0$, $\Delta_g > 0$, such as bilinear query updates along target key $K_{i_\star}$ with correlation margin $\Delta_K$):
\begin{enumerate}
\item \textbf{Linear margin growth and exponential retrieval convergence:} After $T$ steps, the logit margin is at least $\gamma_0 + T\eta\Delta_g$, the attention weight on $i_\star$ satisfies $1 - p_{T,i_\star} \le (n-1)e^{-\gamma_0}e^{-T\eta\Delta_g}$, and for values with $\|v_j - v_{i_\star}\|\le D_V$,
\[
\bigl\|\mathrm{rowOutput}(s_T, m, v) - v_{i_\star}\bigr\| \le (n-1)\,D_V\,e^{-\gamma_0}\,e^{-T\eta\Delta_g}.
\]
\item \textbf{Step complexity and population risk transfer:} Whenever $T\eta\Delta_g \ge \log\bigl(\frac{(n-1)D_V}{\varepsilon}\bigr) - \gamma_0$, the retrieval error is at most $\varepsilon$, and against any oracle regressor using $v_{i_\star}(z)$ with population MSE $R_{\mathrm{oracle}}$, the attention head's population MSE is at most $2 R_{\mathrm{oracle}} + 2\bigl((n-1)D_V e^{-\gamma_0 - T\eta\Delta_g}\bigr)^2$.
\end{enumerate}
\end{lemma}
\begin{proof}
\leanok
Induct on $T$ to show $s_{T,j} \le s_{T,i_\star} - (\gamma_0 + T\eta\Delta_g)$ on unmasked keys $j\ne i_\star$. Apply \cref{lem:attention_logit_stability} with margin $\gamma_T = \gamma_0 + T\eta\Delta_g$ to bound the tail weight and output retrieval error, solve $(n-1)D_V e^{-\gamma_0 - T\eta\Delta_g} \le \varepsilon$ via $\log$, and integrate the pointwise Young inequality $\|\hat{y} - y\|^2 \le 2\|v_{i_\star} - y\|^2 + 2\|\hat{y} - v_{i_\star}\|^2$.
\end{proof}

\begin{lemma}[Multilayer Transformer generalization, in-context excess risk, and safety intervention control]\label{lem:transformer_generalization}
\lean{OpenTheory.TransformerGeneralization.transformerPerturbationEnvelope_smul, OpenTheory.TransformerGeneralization.transformer_stack_param_perturbation_le, OpenTheory.TransformerGeneralization.transformer_depth_scaled_lip_bound, OpenTheory.TransformerGeneralization.transformerParamLoss_lipschitz, OpenTheory.TransformerGeneralization.transformerBall_loss_cover, OpenTheory.TransformerGeneralization.transformer_uniform_deviation_tail, OpenTheory.TransformerGeneralization.transformer_approx_erm_excess_risk_tail, OpenTheory.TransformerGeneralization.icl_end_to_end_excess_risk_bound, OpenTheory.TransformerGeneralization.icl_total_risk_monotone_in_steps, OpenTheory.TransformerGeneralization.icl_end_to_end_target_accuracy, OpenTheory.TransformerGeneralization.safety_intervention_tradeoff_bound, OpenTheory.TransformerGeneralization.orthogonal_circuit_ablation_zero_unsafe_and_risk_bound}
\leanok
\uses{lem:attention_ablation, lem:in_context_retrieval, lem:trajectory_generalization}
For an $L$-layer Transformer residual stack $\mathrm{residualTrajectory}(L, \mathrm{branch}_u, x)$ parameterized by $u\in\mathbb{R}^p$:
\begin{enumerate}
\item \textbf{Depth-scaled parameter-to-output stability and covering generalization:} When layer $t$ has state Lipschitz constant $c_t$ and parameter sensitivity $L_{\mathrm{param},t}$, the stack output is $L_{\mathrm{stack}}$-Lipschitz in $u$ with $L_{\mathrm{stack}} = \mathrm{residualPerturbationEnvelope}(L, c, L_{\mathrm{param}})$, which under $1/L$ residual scaling ($c_t \le c_0/L$, $L_{\mathrm{param},t} \le L_0/L$) satisfies the depth-independent bound $L_{\mathrm{stack}} \le L_0 e^{c_0}$. Over the parameter ball $B_p(R)$, any $L_\ell$-Lipschitz $[0,1]$-loss enjoys uniform empirical-to-population concentration with covering cardinality $N \le (\mathrm{coverAxis}(R,\delta,p)+1)^p$.
\item \textbf{End-to-end In-Context Learning (ICL) excess risk:} Combining $T$ steps of query-key margin alignment ($\gamma_T \ge \gamma_0 + T\eta\Delta_g$) with generalization gap $\varepsilon_{\mathrm{gen}}$ bounds the total population risk by $2 R_{\mathrm{oracle}} + 2\bigl((n-1)D_V e^{-\gamma_0 - T\eta\Delta_g}\bigr)^2 + \varepsilon_{\mathrm{gen}}$, which is monotonically decreasing in $T$ and reaches $2 R_{\mathrm{oracle}} + 2\varepsilon_{\mathrm{ret}}^2 + \varepsilon_{\mathrm{gen}}$ once $T\eta\Delta_g \ge \log\frac{(n-1)D_V}{\varepsilon_{\mathrm{ret}}} - \gamma_0$.
\item \textbf{Mechanistic safety intervention and circuit ablation:} Under layerwise circuit steering $\|\mathrm{branch}_{1,t} - \mathrm{branch}_{2,t}\| \le \varepsilon_{\mathrm{edit},t}$ that suppresses an unsafe linear probe $\|P_{\mathrm{unsafe}}(f_{\mathrm{edit}}(z))\| \le \tau_{\mathrm{safe}}$, benign population MSE is bounded by $2 R_P(f_{\mathrm{clean}}) + 2\,\mathrm{residualPerturbationEnvelope}(L, c, \varepsilon_{\mathrm{edit}})^2$. Moreover, ablating heads $h\notin\mathrm{keep}$ where retained heads satisfy $P_{\mathrm{unsafe}}(\mathrm{head}_h) = 0$ achieves \emph{exact} unsafe-subspace eradication $P_{\mathrm{unsafe}}(\mathrm{MHA}_{\mathrm{keep}}) = 0$ while bounding benign MSE degradation by $2\bigl(\sum_{h\in\mathrm{keep}^c} \|W_{O,h}\|_{\mathrm{op}} B_{V,h}\bigr)^2$.
\end{enumerate}
\end{lemma}
\begin{proof}
\leanok
Factor $\|u-v\|_2$ out of `residualPerturbationEnvelope` by induction on $L$ (`transformerPerturbationEnvelope_smul`) and apply `multilayer_residual_perturbation`. Under $1/L$ depth scaling, bound the recurrence by $L_0(1+c_0/L)^L \le L_0 e^{c_0}$. Combine with `BallCover.euclideanBall_net` and `Covering.uniform_deviation_tail` for the generalization tail. For safety interventions, combine `multilayer_residual_perturbation` and `PredictionGap.sq_norm_le_double` under `P.risk_le_envelope`, and use linearity of $P_{\mathrm{unsafe}}$ over `keep`.
\end{proof}

\begin{lemma}[Convolutional translation equivariance, receptive-field locality, and a comparison of CNN vs.\ MLP covering upper bounds]\label{lem:convolutional_generalization}
\lean{OpenTheory.ConvolutionalGeneralization.cyclicPatch_shift, OpenTheory.ConvolutionalGeneralization.convPreactivation_shift_equivariant, OpenTheory.ConvolutionalGeneralization.convFeatureMap_shift_equivariant, OpenTheory.ConvolutionalGeneralization.convForward_shift_invariant, OpenTheory.ConvolutionalGeneralization.patchEnergy_le_inputEnergy, OpenTheory.ConvolutionalGeneralization.conv_param_diff_le, OpenTheory.ConvolutionalGeneralization.convParamLoss_lipschitz, OpenTheory.ConvolutionalGeneralization.convBall_loss_cover, OpenTheory.ConvolutionalGeneralization.conv_uniform_deviation_tail, OpenTheory.ConvolutionalGeneralization.conv_approx_erm_excess_risk_tail, OpenTheory.ConvolutionalGeneralization.cnn_lip_le_mlp_lip, OpenTheory.ConvolutionalGeneralization.cnn_lip_lt_mlp_lip, OpenTheory.ConvolutionalGeneralization.cnn_covering_strictly_smaller_than_mlp, OpenTheory.ConvolutionalGeneralization.cnn_vs_mlp_sample_complexity_ratio}
\leanok
\uses{lem:trajectory_generalization}
For a $c$-channel convolutional layer with filter width $k \le d$, shared weights $W\in\mathbb{R}^{c\times k}$, activation $\sigma$ ($|\sigma'|\le L$), and readout weights $|\alpha_{\mathrm{ch}}|\le A$ over spatial dimension $d$:
\begin{enumerate}
\item \textbf{Exact translation equivariance and pooling invariance:} For every cyclic spatial shift $s\in\mathrm{Fin}\,d$, the convolutional feature map satisfies $\Phi_{\mathrm{ch}}(\mathrm{shift}_s(x)) = \mathrm{shift}_s(\Phi_{\mathrm{ch}}(x))$, and the global average-pooled predictor is strictly translation invariant: $\mathrm{convForward}(\mathrm{shift}_s(x)) = \mathrm{convForward}(x)$.
\item \textbf{Receptive-field locality vs.\ ambient dimension:} When local patch energy satisfies $\|x_{\mathrm{patch}(p)}\|_2^2 \le E_{\mathrm{patch}} \le E_{\mathrm{global}}$, the parameter-to-loss Lipschitz constant of the CNN is $L_{\mathrm{CNN}} = L_\ell c A L\sqrt{E_{\mathrm{patch}}} \le L_{\mathrm{MLP}} = L_\ell c A L\sqrt{E_{\mathrm{global}}}$ (strictly smaller whenever $E_{\mathrm{patch}} < E_{\mathrm{global}}$).
\item \textbf{Comparison of covering upper bounds ($O(ck)$ vs.\ $O(cd)$):} Whenever $0 < k < d$, the CNN parameter-ball covering bound $(\mathrm{coverAxis}(R,\delta,ck)+1)^{ck}$ is strictly smaller than the dense MLP covering bound $(\mathrm{coverAxis}(R,\delta,cd)+1)^{cd}$, and the ratio of log-covering sample complexities is exactly $k/d < 1$. This compares two \emph{upper} bounds; it is not a lower bound for the MLP and does not prove that CNNs need fewer samples.
\end{enumerate}
\end{lemma}
\begin{proof}
\leanok
Equivariance follows from associativity and commutativity of index addition in $\mathrm{Fin}\,d$, and pooling invariance follows by reindexing the spatial sum along `Equiv.addRight s`. Cauchy--Schwarz on each $k$-dimensional patch bounds the parameter perturbation by $c A L\sqrt{E_{\mathrm{patch}}}\|u-v\|_2$ uniformly across spatial positions. Monotonicity of `coverAxis` in the parameter dimension $ck < cd$ together with strict power growth yields the strict covering and $k/d$ sample complexity separation.
\end{proof}

\chapter{Certified generalization of a trained network}

This chapter is the flagship certificate: a concrete trained network, its exact error count re-computed by the kernel, and a distribution-free bound on its population error. The only unchecked assumption is that the held-out rows are IID and were not used for training.

\begin{lemma}[Finite-code generalization for a trained quantized MLP]\label{lem:quantized_mlp_certificate}
\lean{OpenTheory.QuantizedMLP.uniform_code_tail, OpenTheory.QuantizedMLP.selected_code_tail, OpenTheory.QuantizedMLP.mlp_tail, OpenTheory.QuantizedMLP.numerical_tail_le, OpenTheory.QuantizedMLP.digits_certificate, OpenTheory.QuantizedMLP.empiricalRisk_eq_errorCount, OpenTheory.QuantizedMLPDigits.training_errors, OpenTheory.QuantizedMLPDigits.empirical_error}
\leanok
\uses{lem:reuse_finite_class}
Fix a decoder with $p$ parameters taking $q$ possible values each, a measurable
loss in $[0,1]$, and an arbitrary probability distribution $P$ on observations.
For $n\ge1$ IID observations and any sample-dependent learner $A$ returning a
code, the population risk satisfies
\[
 \Pr_{S\sim P^n}\{R(A(S))>\widehat R_S(A(S))+\varepsilon\}
 \le 2q^p\exp(-2n\varepsilon^2),\qquad \varepsilon>0.
\]
The executable one-hidden-layer ReLU decoder includes both weight matrices and
both biases, so $p=dh+h+hc+c$. For the fixed $16$--$4$--$10$ architecture with
four levels per parameter, $p=118$ and the description length is $236$ bits.
At $n=1400$ and $\varepsilon=1/4$, the failure probability is at most $1/20$.
The saved trained network has exactly $325$ errors on its $1400$ saved training
observations. Its resulting IID-conditional population-error certificate is
$325/1400+1/4=27/56<1/2$.
This is a classical finite-class bound, independent of the optimizer. The
benchmark's IID sampling assumption is external and is not certified by the
saved observations; no guarantee for new writers or distribution shifts is
asserted.
\end{lemma}
\begin{proof}
\leanok
Apply the checked finite-class Hoeffding bound to all $q^p$ codes and take the
one-sided bad-event subset. The simultaneous guarantee permits arbitrary
sample-dependent code selection. Exact exponential bounds discharge the
numerical tail. Integer inference and ordinary kernel reduction check the
saved network's training errors; the empirical-risk bridge identifies that
count with the empirical mean appearing in the probability theorem.
\end{proof}

\begin{theorem}[Catoni's PAC-Bayes bound on a finite hypothesis space]\label{thm:catoni_pac_bayes}
\lean{OpenTheory.PACBayes.catoni_pacBayes, OpenTheory.PACBayes.catoni_grid, OpenTheory.PACBayes.occam_catoni, OpenTheory.PACBayes.exp_moment_le, OpenTheory.PACBayes.expect_le_kl_add_log, OpenTheory.PACBayes.catoni_convexOn, OpenTheory.PACBayes.catoni_le_iff, OpenTheory.PACBayes.kl_prodDist, OpenTheory.PACBayes.kl_uniformOn}
\leanok
Let $D$ be any probability law, let $\ell$ be a measurable loss in $[0,1]$ on a
finite hypothesis space $\mathcal H$, and let $\pi>0$ be a prior fixed
independently of the sample. Write
$\Phi_t(p)=-\log(1-(1-e^{-t})p)$. For every $t\in\mathbb R$ and $\delta>0$,
with probability at least $1-\delta$ over $S\sim D^n$, every posterior $Q$
(possibly depending on $S$) satisfies
\[
 n\,\Phi_t(R(Q)) \le t\,n\,\widehat R_S(Q) + \mathrm{KL}(Q\,\|\,\pi) + \log(1/\delta),
\]
where $R(Q)$ and $\widehat R_S(Q)$ are the population and empirical risks of
the Gibbs predictor. Choosing $t$ from a finite grid $T$ costs $\log|T|$.
A point mass at the output of an arbitrary learner gives the Occam bound with
code length $\log(1/\pi(h))$. For product priors and posteriors on codes, the
relative entropy is the sum of the per-coordinate divergences.
\end{theorem}
\begin{proof}
\leanok
Convexity of $\exp$ dominates each $[0,1]$ loss by a Bernoulli variable, so the
product measure gives $\mathbb E\exp(n\Phi_t(R(h))-t\sum_i\ell(h,z_i))\le1$
for each $h$. Averaging over $\pi$ and Markov's inequality bound the
probability that the $\pi$-average exceeds $1/\delta$. On the complement, the
finite Donsker--Varadhan inequality (Jensen for $\log$) and convexity of
$\Phi_t$ transfer the bound to every $Q$. A union bound handles the grid.
\end{proof}

\begin{lemma}[PAC-Bayes certificate for a trained MLP with a learned prior]\label{lem:pacbayes_mlp_certificate}
\lean{OpenTheory.QuantizedMLP.mlp_pacBayes, OpenTheory.QuantizedMLP.kl_dirac_centre_le, OpenTheory.QuantizedMLP.dirac_certificate, OpenTheory.QuantizedMLP.errorCount_codeOfList, OpenTheory.PACBayesMLPDigits.heldout_errors, OpenTheory.PACBayesMLPDigits.heldout_empirical_error, OpenTheory.PACBayesMLPDigits.heldout_pacBayes, OpenTheory.PACBayesMLPDigits.heldout_certificate, OpenTheory.PACBayesMLPDigits.heldoutSample_fails_iff}
\leanok
\uses{thm:catoni_pac_bayes}
The prior is a product over the $1210$ two-bit parameters of a
$64$--$16$--$10$ ReLU network: each coordinate keeps the level of a reference
network with mass $1-3\cdot2^{-20}$. The reference network was trained, and the
architecture selected, using the $1400$ training rows only. On the $397$
held-out rows the network makes exactly $16$ errors. With $t$ chosen from
$\{1/4,\dots,2\}$ and $\delta=1/20$: with probability at least $95\%$ over IID
held-out observations independent of the training rows, every posterior obeys
Catoni's inequality relative to this prior, and the reference network's
population error is at most $799/10000<8\%$.
The IID assumption and the independence of the two splits are external and
are not certified by the saved observations.
\end{lemma}
\begin{proof}
\leanok
Instantiate Theorem~\ref{thm:catoni_pac_bayes} with the centred product prior
and the Dirac posterior at the reference network, whose relative entropy is at
most $1210\cdot x/(1-x)$ for $x=3\cdot2^{-20}$. Kernel reduction checks the
held-out error count through a linear-time list evaluation of the network.
Rational Taylor bounds on $\exp$ and Mathlib's bounds on $\log 2$ and $e$ invert
$\Phi_t$.
\end{proof}

\begin{theorem}[Exact binomial test-set bound (Clopper--Pearson)]\label{thm:binomial_test_set}
\lean{OpenTheory.BinomialTail.sampleLaw_count_le, OpenTheory.BinomialTail.derivative_sum_bernstein, OpenTheory.BinomialTail.binomCdf_antitoneOn, OpenTheory.BinomialTail.count_testSet_bound, OpenTheory.BinomialTail.count_clopperPearson, OpenTheory.BinomialTail.testSet_bound, OpenTheory.BinomialTail.clopperPearson, OpenTheory.BinomialTail.mem_clopperPearson_event, OpenTheory.BinomialTail.binomCdf_div_eq}
\leanok
Let $D$ be any probability law and $\ell(h,\cdot)\in\{0,1\}$ a measurable loss
of one hypothesis $h$ fixed independently of the sample, with population error
$R$. Write $F_{n,p}(k)=\sum_{i\le k}\binom ni p^i(1-p)^{n-i}$. Under
$S\sim D^n$ the error count $K$ satisfies $\Pr(K\le k)=F_{n,R}(k)$ exactly.
Consequently: if $F_{n,p}(k)\le\delta$ then
$\Pr(K\le k \wedge R>p)\le\delta$; and, uniformly in the observed count,
$\Pr(F_{n,R}(K)\le\delta)\le\delta$, while any realized sample with at most
$k$ errors lies in this event whenever $R>p$ and $F_{n,p}(k)\le\delta$.
\end{theorem}
\begin{proof}
\leanok
The event $K\le k$ is the disjoint union, over sets $A$ of at most $k$
positions, of the measurable rectangles on which the errors fall exactly on
$A$; each has product probability $R^{|A|}(1-R)^{n-|A|}$, and grouping by
$|A|$ gives $F_{n,R}(k)$. In Bernstein polynomials,
$\frac{d}{dp}F_{n,p}(k)=-n\,b_{n-1,k}(p)\le0$ on $[0,1]$, so $F_{n,p}(k)$
decreases in $p$. For the uniform statement let $k^*$ be the largest $k\le n$
with $F_{n,R}(k)\le\delta$; the event is contained in $K\le k^*$.
\end{proof}

\begin{lemma}[Exact binomial certificate for the trained MLP]\label{lem:binomial_mlp_certificate}
\lean{OpenTheory.BinomialMLPDigits.tailNumerator_le, OpenTheory.BinomialMLPDigits.binomCdf_le, OpenTheory.BinomialMLPDigits.heldout_numErrors, OpenTheory.BinomialMLPDigits.heldout_testSet_certificate, OpenTheory.BinomialMLPDigits.heldoutSample_fails_iff, OpenTheory.BinomialMLPDigits.heldout_clopperPearson, OpenTheory.BinomialMLPDigits.heldoutSample_fails_of_risk_gt}
\leanok
\uses{thm:binomial_test_set, lem:pacbayes_mlp_certificate}
For the same saved $64$--$16$--$10$ network and the same $397$ held-out rows
($16$ errors), with $\delta=1/20$: with probability at least $95\%$ over IID
held-out observations independent of the training rows, the population error
is at most $303/5000=6.06\%$ (the PAC-Bayes certificate gives $7.99\%$). The
same number follows from the Clopper--Pearson event, with no error threshold
fixed in advance. The IID assumption and the independence of the two splits
are external and are not certified by the saved observations.
\end{lemma}
\begin{proof}
\leanok
Apply Theorem~\ref{thm:binomial_test_set} with $n=397$, $k=16$,
$p=303/5000$. Writing $F_{397,p}(16)=N/5000^{397}$ with $N$ a natural number,
the inequality $20N\le5000^{397}$ is checked by kernel reduction with
linear-time binomial coefficients.
\end{proof}

\begin{theorem}[Binomial Occam bound for countable hypothesis spaces]\label{thm:occam_binomial}
\lean{OpenTheory.BinomialTail.occam_binomial, OpenTheory.BinomialTail.mem_occamEvent, OpenTheory.BinomialTail.tsum_sigma_mul_le_one, OpenTheory.BinomialTail.tsum_uniformWeight}
\leanok
\uses{thm:binomial_test_set}
Let $\mathcal H$ be countable, each $\ell(h,\cdot)\in\{0,1\}$ measurable, and
$w:\mathcal H\to[0,\infty]$ with $\sum_h w(h)\le1$ (for a prefix-free code,
$w(h)=2^{-L(h)}$). With probability at least $1-\delta$ over $S\sim D^n$, for
every $h$ simultaneously, $F_{n,R_h}(K_h)>\delta\,w(h)$. A realized sample on
which $h$ makes at most $k$ errors, with $F_{n,p}(k)\le\delta\,w(h)$ and
$R_h>p$, lies in the failure event. Weights compose over two-part codes, and
uniform weights on a finite set sum to one.
\end{theorem}
\begin{proof}
\leanok
A countable union bound over Theorem~\ref{thm:binomial_test_set} at level
$\delta\,w(h)$, and $\sum_h\delta\,w(h)\le\delta$.
\end{proof}

\begin{lemma}[Exact dot products by Kronecker substitution]\label{lem:kronecker_dot}
\lean{OpenTheory.Kronecker.ofDigits_conv, OpenTheory.Kronecker.getD_conv, OpenTheory.Kronecker.getD_conv_le, OpenTheory.Kronecker.ofDigits_div_pow_mod, OpenTheory.Kronecker.kronecker_dot, OpenTheory.Kronecker.pack_slots, OpenTheory.Kronecker.slot_lt_of_land_eq_zero, OpenTheory.Kronecker.slots_lt_of_land_eq_zero}
\leanok
For natural lists $c,x$ of length $d$, write
$\mathrm{pack}_B(L)=\sum_k L_k2^{Bk}$. If $x_i\le M$ for all $i$ and
$(\sum_i c_i)M<2^B$, then slot $d-1$ of
$\mathrm{pack}_B(c)\cdot\mathrm{pack}_B(\mathrm{rev}\,x)$ is
$\sum_i c_ix_i$. A number below $2^{Bd}$ is the packing of its first $d$
slots, and if $r\wedge H=0$ for the mask $H$ with bits $m,\dots,B-1$ set in
each of $d$ slots, every one of those slots of $r$ is below $2^m$.
\end{lemma}
\begin{proof}
\leanok
The product is the evaluation at $2^B$ of the carry-free convolution of the
digit lists, whose $k$-th coefficient is $\sum_{i\le k}c_ix_{k-i}\le(\sum c)M$.
Digits below the base are recovered by division and remainder, and the
coefficient at $d-1$ against the reversed list is the dot product. The mask
statement compares bits slot by slot.
\end{proof}

\begin{theorem}[Packed inference equals the reference decoder]\label{thm:packed_mlp_inference}
\lean{OpenTheory.QuantizedMLP.dot_packed, OpenTheory.QuantizedMLP.scoresPacked_eq, OpenTheory.QuantizedMLP.errorCount_packed}
\leanok
\uses{lem:kronecker_dot}
Store each first-layer row of a two-bit MLP as one natural with $B$-bit slots
and each input as one natural holding its pixels in reverse slot order. If the
rows and the all-ones vector are valid packings ($d$ slots, codes below $4$),
inputs are valid packings of pixels in $[0,255]$, and $3\cdot255\,d<2^B$, then
the packed forward pass, with one natural product per hidden unit and the
identity $\sum(c_i-2)x_i=\sum c_ix_i-2\sum x_i$, returns exactly the scores of
the reference decoder, and the error counts agree.
\end{theorem}
\begin{proof}
\leanok
Lemma~\ref{lem:kronecker_dot} for each hidden unit and for the pixel sum; the
reference decoder's row split of the weight list reduces to the packed rows,
and the remaining layers are unchanged.
\end{proof}

\chapter{Design-based certificates from a public seed}

A certifier commits to a pool of $N$ items and to a property $e$ of each item
(for example ``the deployed model's logged answer is wrong''). A public
randomness beacon published after both commitments is expanded into $n$ blocks
of $b$ bits, and block $i$ selects pool index $v_i \bmod N$. The claim is about
the pool, so no distributional assumption is made: the randomness is the seed's.

\begin{theorem}[Soundness of a pool-rate certificate]\label{thm:pool_rate_sound}
\lean{OpenTheory.DesignSampling.rate_sound, OpenTheory.DesignSampling.bad_subset, OpenTheory.DesignSampling.card_deriveIndex_fiber_le, OpenTheory.DesignSampling.card_derive_preimage_le, OpenTheory.DesignSampling.sampleLaw_uniformPool, OpenTheory.DesignSampling.card_clopperPearson_le}
\leanok
\uses{thm:binomial_test_set}
Let $R$ be the fraction of the pool with property $e$, fixed before the seed.
Among all seeds $v\in\{0,\dots,2^b-1\}^n$, those for which some count $k\ge K$
(the observed count) and some $p\ge0$ with $F_{n,p}(k)\le\delta$ satisfy
$p<R$ form a fraction at most $(1+N/2^b)^n\,\delta$.
\end{theorem}
\begin{proof}
\leanok
Each index has at most $\lfloor 2^b/N\rfloor+1$ preimages under reduction
modulo $N$, so a set $A$ of index tuples has at most
$|A|(\lfloor 2^b/N\rfloor+1)^n$ preimage seeds. A bad seed maps into the
Clopper--Pearson event $F_{n,R}(K)\le\delta$. Under $n$ uniform indices this
event is the event of Theorem~\ref{thm:binomial_test_set} for the uniform law
on the pool, whose probability of $e$ is $R$, so it has at most $\delta N^n$
elements. Dividing by $2^{bn}$ gives the bound.
\end{proof}

\begin{corollary}[Selective publication]\label{cor:selective_publication}
\lean{OpenTheory.DesignSampling.selective_sound}
\leanok
\uses{thm:pool_rate_sound}
If an issuer commits $M$ properties $e_1,\dots,e_M$ before the seed and
publishes any one certificate after seeing the sample, the seeds on which the
published claim can be wrong form a fraction at most
$(1+N/2^b)^n\sum_j\delta_j$.
\end{corollary}
\begin{proof}
\leanok
Union bound over the $M$ bad sets of Theorem~\ref{thm:pool_rate_sound}.
\end{proof}

\begin{lemma}[Chernoff witness for the binomial tail]\label{lem:chernoff_witness}
\lean{OpenTheory.ChernoffWitness.binomCdf_le_witness, OpenTheory.ChernoffWitness.binomCdf_le_of_witness}
\leanok
For $p\in[0,1]$ and $t\in(0,1]$, $F_{n,p}(k)\le(1-p+pt)^n/t^k$. With
$p=a/b$, $t=u/v$ and $\delta=1/d$, the integer inequality
$d\,(bv-av+au)^n v^k\le (bv)^n u^k$ implies $F_{n,p}(k)\le\delta$: one
comparison, whatever $k$.
\end{lemma}
\begin{proof}
\leanok
Multiply each term $i\le k$ by $t^{i-k}\ge1$ and add the terms $i>k$; the
binomial theorem sums the result.
\end{proof}

\begin{definition}[Claim template \texttt{pool-rate}]\label{def:pool_rate_claim}
\lean{OpenTheory.PoolRate.Claim, OpenTheory.PoolRate.Claim.of_tail, OpenTheory.PoolRate.Claim.of_witness, OpenTheory.PoolRate.Claim.sound}
\leanok
\uses{thm:pool_rate_sound, lem:chernoff_witness}
A certificate proves that at most $k$ of the $n$ sampled bits are set, that
$p\ge0$ and that $F_{n,p}(k)\le\delta$, either by the exact tail or by a
Chernoff witness. If the bits are those of a property fixed before the seed,
the claim $R\le p$ holds unless the seed lies in the set bounded by
Theorem~\ref{thm:pool_rate_sound}.
\end{definition}

\begin{lemma}[A true decision procedure is a claim]\label{lem:pool_check}
\lean{OpenTheory.FastCheck.poolCheck, OpenTheory.FastCheck.poolCheckWitness, OpenTheory.PoolRate.popTree_eq, OpenTheory.PoolRate.card_bitsOf, OpenTheory.PoolRate.hornerState_eq, OpenTheory.PoolRate.tailFast_eq, OpenTheory.PoolRate.Claim.of_poolCheck, OpenTheory.PoolRate.Claim.of_poolCheckWitness}
\leanok
\uses{def:pool_rate_claim}
The checks \texttt{poolCheck} and \texttt{poolCheckWitness} are \texttt{Bool}-valued functions
in plain natural-number arithmetic, with no import beyond Lean's core: a bit count by halving,
the tail numerator by Horner's rule (linear in $k$), and the Chernoff comparison. If either returns \texttt{true} on $(x,n,k,\dots)$,
the \texttt{pool-rate} claim holds for the bits of $x$. A certificate is therefore checked by
running the function, compiled or in the kernel, without re-proving anything.
\end{lemma}
\begin{proof}
\leanok
The halving count equals the number of set bits among the first $64\cdot2^l$ bits; for $x<2^n$
the bits beyond $n$ vanish. After $k$ Horner steps the state holds $\binom nk$, $a^k$ and
$\sum_{i\le k}\binom ni a^i q^{k-i}$; multiplying by $q^{n-k}$ gives \texttt{tailNumerator}, and for
$k\ge n$ the binomial theorem gives $b^n$.
\end{proof}

\section{A judge and a gold audit}

\begin{lemma}[Rectifier identity, uniform prefixes]\label{lem:rectifier}
\lean{OpenTheory.DesignSampling.poolRate_rectify, OpenTheory.DesignSampling.poolRate_not, OpenTheory.DesignSampling.seedFrac_prefix, OpenTheory.DesignSampling.seedFrac_union_le}
\leanok
For any two properties $T$ (true failure) and $J$ (judged failure) of the pool items,
$R_T = R_J + R_{T\wedge\neg J} - R_{\neg T\wedge J}$ exactly, and $R_{\neg E}=1-R_E$. The first
$m$ blocks of a uniform seed of $n\ge m$ blocks form a uniform seed of $m$ blocks.
\end{lemma}
\begin{proof}
\leanok
Count item by item. Every prefix has exactly $(2^b)^{n-m}$ extensions.
\end{proof}

\begin{theorem}[Exact finite-sample prediction-powered failure rate]\label{thm:judged_rate}
\lean{OpenTheory.JudgedRate.Claim, OpenTheory.JudgedRate.bad, OpenTheory.JudgedRate.bad_frac, OpenTheory.JudgedRate.Claim.sound, OpenTheory.JudgedRate.Claim.sound_alarms}
\leanok
\uses{thm:pool_rate_sound, lem:rectifier, def:pool_rate_claim}
Let a judge's verdicts be committed for every pool item. From a beacon sample of $n$ items and
gold labels on its first $m$, certificates $R_J\le p_J$ (judge failures, all $n$),
$R_{T\wedge\neg J}\le p_A$ (missed failures, first $m$) and optionally $R_{\neg(\neg T\wedge J)}\le p_B$
(first $m$) give $R_T\le p_J+p_A$, or $R_T\le p_J+p_A+p_B-1$, except on a set of seeds of fraction
at most $(1+N/2^b)^n\delta_J+(1+N/2^b)^m(\delta_A+\delta_B)$.
\end{theorem}
\begin{proof}
\leanok
Apply Theorem~\ref{thm:pool_rate_sound} to each part, on the seed or its prefix, and take a union
bound; outside the bad set the rectifier identity combines the three bounds.
\end{proof}

\chapter{Gradient-descent plateau and singular learning}

\begin{theorem}[The gradient-descent plateau: exact mode dynamics]\label{thm:plateau_modes}
\lean{OpenTheory.Plateau.mode_zero, OpenTheory.Plateau.hasDerivAt_mode, OpenTheory.Plateau.sub_mode_le, OpenTheory.Plateau.mode_le, OpenTheory.Plateau.signal_learned, OpenTheory.Plateau.noise_off, OpenTheory.Plateau.plateau}
\leanok
For $0<\varepsilon^2<\sigma$, $w(t)=\sigma/(1+(\sigma/\varepsilon^2-1)e^{-2\sigma t})$
satisfies $w(0)=\varepsilon^2$ and $\dot w=2w(\sigma-w)$; this is the strength of one
singular mode of a two-layer linear network trained by gradient flow on whitened
data from a balanced initialization of scale $\varepsilon$ aligned with the data.
For finitely many modes $\sigma_i>0$ and $c_1<c_2$ with $\sigma_ic_1>1$ for signal
modes and $\sigma_ic_2<1$ for noise modes, for every $\delta>0$ and all small
enough $\varepsilon>0$, uniformly for $t\in[c_1\log(1/\varepsilon),c_2\log(1/\varepsilon)]$,
every signal mode is within $\delta$ of $\sigma_i$ and every noise mode is within
$\delta$ of $0$.
\end{theorem}
\begin{proof}
\leanok
Explicit bounds $0\le\sigma-w(t)\le\sigma(\sigma/\varepsilon^2)e^{-2\sigma t}$ and
$0\le w(t)\le\sigma/((\sigma/\varepsilon^2-1)e^{-2\sigma t})$; at
$t=c\log(1/\varepsilon)$, $e^{-2\sigma t}=\varepsilon^{2\sigma c}$, and the bounds are
monotone in $t$ on the window.
\end{proof}

\begin{theorem}[Mode decoupling for a two-layer linear network]\label{thm:plateau_flow}
\lean{OpenTheory.Plateau.unique_of_contDiff, OpenTheory.Plateau.gradField_eq_matrix, OpenTheory.Plateau.gradField_contDiff, OpenTheory.Plateau.diagSol_solves, OpenTheory.Plateau.flow_exists, OpenTheory.Plateau.flow_decouples, OpenTheory.Plateau.flow_plateau, OpenTheory.Plateau.gradField_rotate, OpenTheory.Plateau.rotate_solves}
\leanok
\uses{thm:plateau_modes}
Gradient flow of $\tfrac12\|BA-C\|^2$ is $\dot A=B^\top(C-BA)$,
$\dot B=(C-BA)A^\top$. For diagonal data $C=\operatorname{diag}(\sigma_i)$ with
$\varepsilon^2<\sigma_i$, width $H\le\min(M,N)$, and initialization
$A(0)=\varepsilon I$, $B(0)=\varepsilon I$, every solution on $[0,T]$ satisfies
$(BA)(t)_{ij}=\delta_{ij}\,[i<H]\,w_i(t)$ with $w_i$ the logistic mode of
Theorem~\ref{thm:plateau_modes}; such a solution exists. Consequently, on the
window $[c_1\log(1/\varepsilon),c_2\log(1/\varepsilon)]$, every entry of $BA$ is
within $\delta$ of the best fit keeping exactly the signal modes, for all small
$\varepsilon$. The field is equivariant under $(A,B)\mapsto(AV,U^\top B)$ for
orthogonal $U,V$, so with an SVD $C=U\Sigma V^\top$ the general case reduces to
the diagonal one.
\end{theorem}
\begin{proof}
\leanok
The diagonal candidate with $a_k=b_k=\sqrt{w_k}$ solves the flow entrywise;
uniqueness holds because the field is $C^1$, both solutions are bounded on the
compact interval, and a $C^1$ field is Lipschitz on a closed ball.
\end{proof}

\begin{theorem}[Bayes second moment for a scale-invariant prior]\label{thm:bayes_scale_invariant}
\lean{OpenTheory.BayesToy.integrable_abs_rpow_mul_exp, OpenTheory.BayesToy.hasDerivAt_q, OpenTheory.BayesToy.hasDerivAt_q1, OpenTheory.BayesToy.integral_Fode', OpenTheory.BayesToy.ode, OpenTheory.BayesToy.q1_sq_le, OpenTheory.BayesToy.abs_L1_le, OpenTheory.BayesToy.gaussian_ibp, OpenTheory.BayesToy.moment_identity, OpenTheory.BayesToy.postMean_eq, OpenTheory.BayesToy.bayes_second_moment}
\leanok
Let $\alpha>0$, observe $z\sim N(u,1)$ with the improper prior $|u|^{\alpha-1}du$, and let
$m(z)$ be the posterior mean. Under the truth $u=0$, i.e.\ $z\sim N(0,1)$,
$\mathbb E[m(z)^2]=\alpha$. For the toy models $y=ab^kx$ ($\alpha=1/k$) this is the
Bayes generalization coefficient $\alpha/2=1/(2k)$, the learning coefficient $\lambda$.
\end{theorem}
\begin{proof}
\leanok
With $q=\int|u|^{\alpha-1}e^{-(z-u)^2/2}du$, differentiation under the integral sign
and integration by parts in $u$ (split at $0$) give $q''+zq'=(\alpha-1)q$. With
$L=\log q$, $m=z+L'$ and $L'^2=(\alpha-1)-zL'-L''$; a posterior Cauchy--Schwarz
inequality and the equation give $|L'|\le|z|+\sqrt\alpha$, and Gaussian integration by
parts gives $\mathbb E[L'']=\mathbb E[zL']$ and $\mathbb E[z^2]=1$, hence
$\mathbb E[m^2]=\alpha$.
\end{proof}

\begin{lemma}[Bayes pays more than the true dimension in reduced-rank regression]\label{lem:rrr_rlct_gain}
\lean{OpenTheory.RRR.rlct8, OpenTheory.RRR.rlct8_gain}
\leanok
With the learning coefficient $\lambda(M,N,H,r)$ of reduced-rank regression
(Aoyagi--Watanabe 2005, cited), whenever $r<\min(M,N)$ and $r<H$,
$\lambda\ge r(M+N-r)/2+1/2$.
\end{lemma}
\begin{proof}
\leanok
Three branches of the formula reduce to $(H-r)(M-r)/2$, $(H-r)(N-r)/2$ and
$(M-r)(N-r)/2$. In the main branch, $8\lambda-4r(M+N-r)=2g(S-h)+g^2-D^2+\pi$ with
$g=H-r$, $h=H+r$, $S=M+N$, $D=M-N$ and parity $\pi$; case analysis on $S-h$.
\end{proof}

\begin{lemma}[Second-order perturbation of the truncated SVD]\label{lem:truncation_second_order}
\lean{OpenTheory.HilbertSchmidt.hs, OpenTheory.HilbertSchmidt.hs_comp_left_le, OpenTheory.HilbertSchmidt.hs_comp_right_le, OpenTheory.HilbertSchmidt.opNorm_le_hs, OpenTheory.RankTruncation.projL_comp, OpenTheory.Perturbation.opn_gramL_sub_le, OpenTheory.Perturbation.second_order_bound}
\leanok
Let $W,C:E\to F$, $D=C-W$, suppose every nonzero eigenvalue of $WW^\top$ and $W^\top W$ is
at least $s_2$, and let $0<\tau<s_2$, $\delta=s_2-\tau$. With $P,Q$ the spectral
projections of $WW^\top,W^\top W$ onto $[s_2,\infty)$, $P_T(D)=D-P^\perp DQ^\perp$,
$h_L=\|CC^\top-WW^\top\|$, $h_R=\|C^\top C-W^\top W\|$,
\[\|P_{[\tau,\infty)}(CC^\top)C-W-P_T(D)\|_{HS}\le\tfrac{h_Lh_R}{\delta^2}\|W\|_{HS}
+\big(\tfrac{h_L+h_R}{\delta}+\tfrac{h_L+h_R}{\tau}\big)\|D\|_{HS}.\]
\end{lemma}
\begin{proof}
\leanok
Exact identity $\widehat W-W-P_T(D)=-[P_C^\perp P_W WQ_WQ_C^\perp+(P_C^\perp-P_W^\perp)DQ_C^\perp
+P_W^\perp D(Q_C^\perp-Q_W^\perp)]$ from the intertwining $P_C C=CQ_C$; each projection
product is bounded by Davis--Kahan (vendored SLT library); Hilbert--Schmidt norms by
operator norm times Hilbert--Schmidt norm.
\end{proof}

\begin{lemma}[Gaussian trace identity]\label{lem:gaussian_trace}
\lean{OpenTheory.GaussianQuadratic.integral_inner_mul_inner, OpenTheory.GaussianQuadratic.integral_norm_sq_map, OpenTheory.GaussianQuadratic.integrable_one_add_norm_pow}
\leanok
For $x$ standard Gaussian on a finite-dimensional space with orthonormal basis $(g_a)$ and
$L$ linear, $\mathbb E\|Lx\|^2=\sum_a\|Lg_a\|^2$; all moments of $\|x\|$ are finite.
\end{lemma}
\begin{proof}
\leanok
The covariance of Mathlib's \texttt{stdGaussian} is the inner product; Fernique for moments.
\end{proof}

\begin{theorem}[Asymptotic risk of the truncated SVD is the manifold dimension]\label{thm:truncation_risk}
\uses{lem:truncation_second_order, lem:gaussian_trace}
\lean{OpenTheory.TruncationRisk.tendsto_risk, OpenTheory.TruncationRank.rankL_eq_rankR, OpenTheory.TruncationRank.sum_hs_sq_tangent_noise, OpenTheory.TruncationRank.tendsto_risk_iid}
\leanok
Let $W:\mathbb R^M\to\mathbb R^N$ have its nonzero squared singular values $\ge s_2$, let
$r$ be their number, $0<\tau<s_2$, and $\widehat W_n$ the singular-value truncation at
$\tau$ of $W+n^{-1/2}Z$ with $Z$ an $N\times M$ matrix of i.i.d.\ $N(0,1)$ entries. Then
\[n\,\mathbb E\|\widehat W_n-W\|_F^2\longrightarrow r(M+N-r),\]
the dimension of the manifold of rank-$r$ matrices (the expectation is a lower Lebesgue
integral, so no measurability of $\widehat W_n$ is assumed). The counts of eigenvalues
$\ge s_2$ of $WW^\top$ and $W^\top W$ agree.
\end{theorem}
\begin{proof}
\leanok
Lemma~\ref{lem:truncation_second_order} with $D=n^{-1/2}Z$ gives
$|\sqrt n(\widehat W_n-W)-P_T(Z)|_{HS}\le K(\|Z\|)/\sqrt n$ with $K$ a polynomial;
squeeze between two integrals with an integrable envelope $c(1+\|Z\|)^8$; the limit
$\mathbb E\|P_TZ\|^2=\sum_{k,l}(1-\|P^\perp f_k\|^2\|Q^\perp e_l\|^2)=NM-(N-r)(M-r)$
by Lemma~\ref{lem:gaussian_trace} and the trace of spectral projections.
\end{proof}

\begin{lemma}[Small noise fixes the spectral picture]\label{lem:truncation_good_event}
\lean{OpenTheory.TruncationEvent.le_inner_of_mem', OpenTheory.TruncationEvent.inner_le_of_mem, OpenTheory.TruncationEvent.eigen_split, OpenTheory.TruncationEvent.count_eq_of_close, OpenTheory.TruncationEvent.eigenvalue_split, OpenTheory.TruncationEvent.card_le_count}
\leanok
If $\|CC^\top-WW^\top\|\le h$ and the nonzero eigenvalues of $WW^\top$ are $\ge s_2$, every
eigenvalue of $CC^\top$ is $\le h$ or $\ge s_2-h$ (Weyl); for $h<\tau<s_2-h$ exactly
$\operatorname{rank}W$ of them are $\ge\tau$.
\end{lemma}
\begin{proof}
\leanok
Weyl's inequality (vendored SLT), Rayleigh bounds on spectral subspaces, and the dimension
of a spectral subspace.
\end{proof}

\begin{theorem}[The gradient-descent plateau is the truncated SVD and pays the true dimension]\label{thm:gd_plateau_end_to_end}
\uses{thm:truncation_risk, lem:truncation_good_event, lem:rrr_rlct_gain}
\lean{OpenTheory.Plateau.trunc_eq_svd, OpenTheory.Plateau.rotated_plateau, OpenTheory.Plateau.plateau_eq_trunc, OpenTheory.Plateau.gd_plateau_end_to_end}
\leanok
Two-layer linear network, whitened design, data $\widehat C_n=W_0+n^{-1/2}Z$ with
i.i.d.\ $N(0,1)$ noise, $W_0$ with $r$ nonzero squared singular values all $\ge s_2$,
$r\le H\le\min(M,N)$. Off events of probability $\to0$, for every SVD of $\widehat C_n$
with positive singular values, gradient flow from the aligned start $\varepsilon$ reaches
at time $\sqrt{2/s_2}\log(1/\varepsilon)$, as $\varepsilon\to0$, the truncated SVD
$\widehat W_n$; $n\,\mathbb E\|\widehat W_n-W_0\|_F^2\to r(M+N-r)$; and when $r<\min(M,N,H)$
the Aoyagi--Watanabe coefficient satisfies $\lambda\ge r(M+N-r)/2+1/2$.
\end{theorem}
\begin{proof}
\leanok
In SVD coordinates the spectral projection keeps exactly the columns with $\sigma_k^2\ge\tau$,
so $\mathrm{trunc}_\tau\widehat C=U\,\mathrm{diag}(\sigma\mathbf 1_{\sigma^2\ge\tau})V^\top$;
rotate the plateau theorem by $(U,V)$; on the good event the signal modes satisfy
$\sigma_k^2\ge3s_2/4$ and the noise modes $\sigma_k^2\le s_2/4$, and no mode beyond $H$ passes
the threshold (too many orthonormal vectors in an $r$-dimensional subspace). Chebyshev for
the good event; Theorem~\ref{thm:truncation_risk} for the risk; Lemma~\ref{lem:rrr_rlct_gain}.
\end{proof}

\begin{theorem}[A loss-threshold stopping rule provably lands on the plateau]\label{thm:stopping_rule}
\uses{thm:gd_plateau_end_to_end}
\lean{OpenTheory.Plateau.trainLoss, OpenTheory.Plateau.hitTime, OpenTheory.Plateau.trainLoss_rotate, OpenTheory.Plateau.rotated_plateau_window, OpenTheory.Plateau.plateauMatrix_eq, OpenTheory.Plateau.stopping_rule, OpenTheory.Plateau.good_event, OpenTheory.Plateau.stopping_rule_end_to_end}
\leanok
Let $\hat t$ be the first time the training loss $\tfrac12\|\widehat C-BA\|_F^2$ is at most $\rho$.
For any fixed $0<\rho\le s_2/4$, off events of probability $\to0$, gradient flow from the
aligned start $\varepsilon$ stopped at $\tfrac32\hat t$ converges, as $\varepsilon\to0$, to the
truncated SVD, whose risk is $r(M+N-r)$ (Theorem~\ref{thm:gd_plateau_end_to_end}).
\end{theorem}
\begin{proof}
\leanok
The loss is rotation invariant, so it can be computed in SVD coordinates, where the network
product is explicit. Before $0.9\log(1/\varepsilon)/\sigma_*$ the weakest signal mode is still off
and the loss exceeds $\tfrac12(7\sigma_*/8)^2>\rho$. At $1.1\log(1/\varepsilon)/\sigma_*$ the loss is
at most $\tfrac N2(\sqrt h+\eta)^2<\rho$. Hence $\tfrac32\hat t\in[1.35,1.65]\log(1/\varepsilon)/\sigma_*$,
which lies inside the plateau window because the noise singular values are below $\sigma_*/\sqrt7$.
\end{proof}

\begin{lemma}[The matrix Riccati equation in closed form]\label{lem:riccati_closed_form}
\lean{OpenTheory.Riccati.expS, OpenTheory.Riccati.phi, OpenTheory.Riccati.Gop, OpenTheory.Riccati.Pc, OpenTheory.Riccati.isUnit_Gop, OpenTheory.Riccati.hasDerivAt_Pc, OpenTheory.Riccati.Pc_zero, OpenTheory.Riccati.eq_Pc}
\leanok
For symmetric $S=\sum_a\lambda_aw_aw_a^\top$ and $Z_0$, $P(t)=Y G^{-1}Y^\top$ with $Y=e^{tS}Z_0$,
$G=I+\sum_a\varphi(\lambda_a,t)g_ag_a^\top$, $g_a=Z_0^\top w_a$, $\varphi(\lambda,t)=\int_0^te^{2\lambda s}ds$,
is the unique solution of $\dot P=SP+PS-P^2$, $P(0)=Z_0Z_0^\top$.
\end{lemma}
\begin{proof}
\leanok
$\dot Y=SY$, $\dot G=Y^\top Y$ and $\tfrac{d}{dt}G^{-1}=-G^{-1}\dot GG^{-1}$ (derivative of the ring
inverse); uniqueness from the smooth vector field.
\end{proof}

\begin{lemma}[Riccati entries on the plateau window]\label{lem:riccati_window}
\uses{lem:riccati_closed_form}
\lean{OpenTheory.Riccati.inner_Q_sq_le, OpenTheory.Riccati.inner_Q_signal_le, OpenTheory.Riccati.dual_identity, OpenTheory.Riccati.amp_small, OpenTheory.Riccati.weight_small, OpenTheory.Riccati.signal_ratio, OpenTheory.Riccati.inner_w_Pc, OpenTheory.Riccati.Pc_entries_limit}
\leanok
With $Z_0=\varepsilon Z_1$, if the signal directions $Z_1^\top w_a$ ($\lambda_ac_1>1$) admit dual
vectors and all other modes have $\lambda_ac_2<1$, then uniformly for
$t\in[c_1,c_2]\log(1/\varepsilon)$, $\langle w_a,P(t)w_b\rangle\to2\lambda_a\mathbf 1_{a=b\text{ signal}}$.
\end{lemma}
\begin{proof}
\leanok
$0\le Q=G^{-1}\le I$, Cauchy--Schwarz for $Q$, the bound $\langle cg,Qcg\rangle\le c^2/\varphi$, and
the dual identity $\langle h,d\rangle=\langle Qh,d\rangle+\sum_c\varphi_c\langle g_c,Qh\rangle\langle g_c,d\rangle$.
\end{proof}

\begin{theorem}[The plateau from a balanced, non-aligned initialization]\label{thm:balanced_plateau}
\uses{lem:riccati_window, thm:gd_plateau_end_to_end}
\lean{OpenTheory.Plateau.gap_const, OpenTheory.Plateau.Zmat_grad, OpenTheory.Plateau.hasDerivAt_Pop, OpenTheory.Plateau.wD, OpenTheory.Plateau.Sop_wD, OpenTheory.Plateau.target_expand, OpenTheory.Plateau.balanced_plateau_diag, OpenTheory.Plateau.balanced_plateau_eq_trunc}
\leanok
Gradient flow of the two-layer linear network from any balanced start $\varepsilon(A_0,B_0)$
($A_0A_0^\top=B_0^\top B_0$) whose signal directions $A_0v_k+B_0^\top u_k$ are linearly independent
(dual vectors) reaches, on the good event and at time $\sqrt{2/s_2}\log(1/\varepsilon)$ as
$\varepsilon\to0$, the truncated SVD. No alignment of the start with the data is assumed.
\end{theorem}
\begin{proof}
\leanok
Balance is conserved, so $Z=[A^\top;B]$ satisfies $\dot Z=SZ-\tfrac12ZZ^\top Z$ with $S$ the dilation of
$C$, and $P=ZZ^\top$ solves the Riccati equation; $BA$ is its off-diagonal block. The dilation of
diagonal data has an explicit eigenbasis; Lemma~\ref{lem:riccati_window} gives the limit; rotate by
the SVD and identify with the truncated SVD.
\end{proof}

\chapter{Resolution of singularities and the learning coefficient}

The resolution theorem below is Resolution's formalization of Hironaka's theorem
(\texttt{resolutionorg/hironaka}, tag \texttt{v4.34.0}), a pinned Lake dependency whose proof is
checked in this build. The rest of the chapter derives from it the existence and rationality of the
learning coefficient of every analytic model (Watanabe 2009, Theorem~6.1 in its deterministic form).

\begin{definition}[Real log canonical threshold]\label{def:rlct}
\lean{OpenTheory.Singular.zetaIntegral, OpenTheory.Singular.IsRLCT, OpenTheory.Singular.IsRLCT.unique}
\leanok
For $K,\varphi$ on a measure space, $\lambda$ is the real log canonical threshold of $K$ against
$\varphi$ if, for every $s\ge0$, $\int K^{-s}\varphi<\infty$ exactly when $s<\lambda$ (the integrand
is $+\infty$ where $K=0$ and $s>0$). A positive threshold is unique.
\end{definition}

\begin{theorem}[Resolution in the form used for Bayesian integrals]\label{thm:hironaka_monomial_bounds}
\lean{OpenTheory.Singular.chartMonomial, OpenTheory.Singular.IsMonomialAt.exists_abs_bounds, OpenTheory.Singular.MonomialBounds, OpenTheory.Singular.exists_resolution_monomial_bounds}
\leanok
Let $K$ be real-analytic on an open $U\subseteq\mathbb R^n$, not identically zero near any point.
There is a proper real-analytic $\pi:M\to U$ from an $n$-dimensional analytic manifold, bijective
over $\{K\neq0\}$, such that near every $p\in M$, in the chart coordinates $y$ at $p$,
$c|y^{\alpha}|\le|K\circ\pi|\le C|y^{\alpha}|$ and $c|y^{\beta}|\le|\det D\pi|\le C|y^{\beta}|$ for
some exponents $\alpha,\beta$ and $0<c\le C$; the preimage of every compact set is covered by
finitely many such neighbourhoods.
\end{theorem}
\begin{proof}
\leanok
Atiyah's resolution theorem with the Jacobian clause of Bierstone--Milman,
\texttt{exists\_proper\_analytic\_monomialization} upstream, gives $K\circ\pi$ and $\det D\pi$ as
monomials times analytic units near every point; a unit is bounded between two positive constants
near the point, and $\pi$ is proper.
\end{proof}

\begin{lemma}[Tonelli for products of one-variable functions]\label{lem:lintegral_fin_prod}
\lean{OpenTheory.Singular.lintegral_fin_prod_eq_prod}
\leanok
For $\sigma$-finite $\mu_i$ on $\mathbb R$ and measurable $f_i\ge0$,
$\int\prod_if_i(x_i)\,d(\bigotimes_i\mu_i)=\prod_i\int f_i\,d\mu_i$.
\end{lemma}
\begin{proof}
\leanok
Induction on the number of variables, splitting off the first coordinate.
\end{proof}

\begin{theorem}[The learning coefficient of a normal-crossing model]\label{thm:normal_crossing_rlct}
\lean{OpenTheory.Singular.cube, OpenTheory.Singular.monomialLoss, OpenTheory.Singular.monomialPrior, OpenTheory.Singular.lintegral_Ioo_rpow_lt_top_iff, OpenTheory.Singular.lintegral_Ioo_rpow_pos, OpenTheory.Singular.zetaIntegrand_eq_prod, OpenTheory.Singular.zetaIntegral_monomial, OpenTheory.Singular.isRLCT_monomial}
\leanok
\uses{def:rlct, lem:lintegral_fin_prod}
On $(0,1)^d$ with $K(y)=\prod_iy_i^{2k_i}$ and $\varphi(y)=\prod_iy_i^{h_i}$, the real log canonical
threshold is $\lambda=\min_{k_i>0}(h_i+1)/(2k_i)$.
\end{theorem}
\begin{proof}
\leanok
$K^{-s}\varphi=\prod_iy_i^{h_i-2k_is}$ factorizes; each factor $\int_0^1y^{a}dy$ is positive and
finite iff $a>-1$.
\end{proof}

\begin{theorem}[Zeta function of a normal-crossing model]\label{thm:normal_crossing_zeta}
\lean{OpenTheory.Singular.monomialLoss_rpow_mul, OpenTheory.Singular.integral_Ioo_rpow, OpenTheory.Singular.zeta_monomial}
\leanok
\uses{thm:normal_crossing_rlct}
For real $z$ with $2k_iz+h_i>-1$ for all $i$,
$\int_{(0,1)^d}K^z\varphi=\prod_i1/(2k_iz+h_i+1)$; its poles $-(h_i+1)/(2k_i)$ have largest
element $-\lambda$.
\end{theorem}
\begin{proof}
\leanok
Fubini and $\int_0^1y^{a}dy=1/(a+1)$.
\end{proof}

\begin{theorem}[Free energy of a normal-crossing model]\label{thm:normal_crossing_free_energy}
\lean{OpenTheory.Singular.lintegral_Ioo_pow, OpenTheory.Singular.measurable_monomialLoss, OpenTheory.Singular.measurable_monomialPrior, OpenTheory.Singular.partition_lower, OpenTheory.Singular.partition, OpenTheory.Singular.freeEnergy_le}
\leanok
\uses{lem:lintegral_fin_prod}
If $k_{i_0}>0$ and $n\ge1$, $Z(n)=\int_{(0,1)^d}e^{-nK}\varphi\ge
e^{-1}n^{-(h_{i_0}+1)/(2k_{i_0})}\prod_i1/(h_i+1)$, so
$-\log Z(n)\le\frac{h_{i_0}+1}{2k_{i_0}}\log n+1+\sum_i\log(h_i+1)$.
\end{theorem}
\begin{proof}
\leanok
On the box where $y_{i_0}<n^{-1/(2k_{i_0})}$, $nK\le1$; integrate $\varphi$ over the box.
\end{proof}

\begin{corollary}[The toy model $y=ab^kx$]\label{cor:toy_rlct}
\lean{OpenTheory.Singular.isRLCT_toy}
\leanok
\uses{thm:normal_crossing_rlct}
$K(a,b)=a^2b^{2k}$ with a uniform prior on $(0,1)^2$ has $\lambda=1/(2k)$ for $k\ge1$, the Bayes
coefficient $\alpha/2$ of Theorem~\ref{thm:bayes_scale_invariant} at $\alpha=1/k$.
\end{corollary}
\begin{proof}
\leanok
$\min(1/2,1/(2k))=1/(2k)$.
\end{proof}

\begin{lemma}[Monomials on coordinate boxes]\label{lem:monomial_box}
\lean{OpenTheory.Singular.intervalIntegrable_abs_rpow, OpenTheory.Singular.lintegral_Ioo_abs_rpow_lt_top_iff, OpenTheory.Singular.lintegral_Ioo_abs_rpow_lt_top_of_lt, OpenTheory.Singular.lintegral_Ioo_abs_rpow_pos, OpenTheory.Singular.box, OpenTheory.Singular.box_subset_ball, OpenTheory.Singular.lintegral_box_prod, OpenTheory.Singular.ae_coord_ne_zero, OpenTheory.Singular.lintegral_box_monomial_lt_top_iff, OpenTheory.Singular.lintegral_finset_biUnion_le}
\leanok
\uses{lem:lintegral_fin_prod}
Let $B\subseteq\mathbb R^n$ be the box of half-width $r>0$ around $c$, with every $c_i$ either $0$ or
of absolute value $>r$. Then $\int_B\prod_i|y_i|^{a_i}dy<\infty$ iff $a_i>-1$ for every $i$ with
$c_i=0$. Almost every point of $\mathbb R^n$ has no zero coordinate.
\end{lemma}
\begin{proof}
\leanok
Fubini on the box; $\int_{-r}^{r}|t|^adt<\infty$ iff $a>-1$, and $|t|^a$ is continuous on an
interval avoiding $0$; every factor is positive.
\end{proof}

\begin{theorem}[Every analytic model has a rational learning coefficient]\label{thm:analytic_rlct}
\lean{OpenTheory.Singular.chartMap, OpenTheory.Singular.LocalBox, OpenTheory.Singular.localBox_nonempty, OpenTheory.Singular.surjective_of_resolution, OpenTheory.Singular.abs_monomial_mul_rpow, OpenTheory.Singular.exists_isRLCT}
\leanok
\uses{def:rlct, thm:hironaka_monomial_bounds, lem:monomial_box}
Let $K\ge0$ be real-analytic on an open $U\subseteq\mathbb R^n$, not identically zero near any point,
and let $\varphi$ be continuous with compact support in $U$, positive at every zero of $K$ in its
support, with at least one such zero. Then $K$ has a positive rational real log canonical threshold
against $\varphi$.
\end{theorem}
\begin{proof}
\leanok
Resolve $K$ by $\pi:M\to U$. Around every $p\in M$ take a box in the chart at $p$ on which
$K\circ\pi\asymp|y^{\alpha}|$, $|\det D\pi|\asymp|y^{\beta}|$, every centre coordinate is $0$ or
exceeds the half-width, and $\varphi\circ\pi\ge\varphi(\pi(p))/2$. $\pi$ is proper, so finitely many
boxes cover $\pi^{-1}(\operatorname{supp}\varphi)$, and $\pi$ is onto since its range is closed and
contains the dense set $\{K\neq0\}$. Let $\lambda$ be the least $(\beta_i+1)/\alpha_i$ over boxes and
coordinates with vanishing centre and $\alpha_i>0$; a zero of $K$ in the support provides one. By the
change of variables formula each piece of $\int K^{-s}\varphi$ is between constant multiples of
$\int_B\prod|y_i|^{\beta_i-s\alpha_i}$, finite for $s<\lambda$ by Lemma~\ref{lem:monomial_box}. The
zero set of $K$ in the support is null, covered by images of coordinate hyperplanes. At the
minimizing box $\varphi$ is bounded below, and the piece diverges for $s\ge\lambda$.
\end{proof}

\chapter{Exact integer networks and region certificates}

A network has two forms. The proof form is a TorchLean chain read over $\mathbb R$; the check form
is an integer runner that the kernel evaluates in milliseconds. Each primitive carries an agreement
proof, and every guarantee below is proved once for all networks built from such primitives, by
induction on the composition. Primitives are vector maps $[m]\to[n]$ with integer parameters;
matrix-shaped primitives such as attention are not covered. These are formal results: no method in
the catalog uses them yet, so they carry no assurance level.

\begin{definition}[Exact primitive and network]\label{def:exact_net}
\lean{OpenTheory.Exact.Prim, OpenTheory.Exact.Net, OpenTheory.Exact.Net.run, OpenTheory.Exact.Net.eval, OpenTheory.Exact.Net.Rules}
\leanok
An exact primitive $[m]\to[n]$ is a TorchLean primitive together with integer parameters, a Boolean
shape check on them, an integer runner on lists, and a proof that on valid parameters and inputs of
length $m$ the runner computes the TorchLean semantics over $\mathbb R$. A network is a primitive or
the sequential composition of two networks; its runner and its real semantics $\mathrm{eval}$ are
the compositions of those of its primitives. For a contract $C$ on primitives, $\mathrm{Rules}_C(g)$
holds one instance of $C$ per primitive of $g$; it describes the class of networks built from
primitives that meet $C$.
\end{definition}

\begin{theorem}[Agreement for every network]\label{thm:exact_agrees}
\lean{OpenTheory.Exact.Net.agrees}
\leanok
\uses{def:exact_net}
For every network $g$, valid parameters $\theta$ and integer input $x$ of length $m$, the integer
runner and the TorchLean model agree: $\mathrm{vec}(g.\mathrm{run}\,\theta\,x) =
g.\mathrm{eval}\,\theta\,(\mathrm{vec}\,x)$.
\end{theorem}
\begin{proof}
\leanok
Induction on the composition. A primitive is its own agreement proof; for $f$ followed by $g$,
the parameter pack splits back into the two packs and the two hypotheses compose.
\end{proof}

\begin{definition}[Integer boxes and interval rules]\label{def:interval_rule}
\lean{OpenTheory.Exact.InBox, OpenTheory.Exact.IntervalRule, OpenTheory.Exact.Net.box}
\leanok
\uses{def:exact_net}
A box is a list of integer pairs $(\ell_i,u_i)$; a real vector $x$ lies in it when the lengths match
and $\ell_i\le x_i\le u_i$ for every $i$. An interval rule for a primitive maps input boxes to output
boxes and is sound: every real input in the input box is sent by the TorchLean semantics into the
output box. Interval propagation through a network composes the rules of its primitives.
\end{definition}

\begin{lemma}[Exact linear layer and ReLU]\label{lem:exact_primitives}
\lean{OpenTheory.Exact.linear_agrees, OpenTheory.Exact.linear_box_sound, OpenTheory.Exact.relu_agrees, OpenTheory.Exact.relu_box_sound}
\leanok
\uses{def:exact_net, def:interval_rule}
TorchLean's linear layer $x\mapsto Wx+b$ with integer $W,b$ and TorchLean's ReLU are exact
primitives with interval rules. For the linear layer the output bounds of row $i$ are
$\sum_k \min(w_{ik}\ell_k, w_{ik}u_k)+b_i$ and $\sum_k \max(w_{ik}\ell_k, w_{ik}u_k)+b_i$; for ReLU
they are $\max(\ell_i,0)$ and $\max(u_i,0)$.
\end{lemma}
\begin{proof}
\leanok
Coordinatewise: TorchLean's matrix--vector product is a sum of products, the integer dot product
is the same sum, and each weighted term $w x_k$ is bounded by the endpoint chosen by the sign of $w$.
ReLU is monotone.
\end{proof}

\begin{theorem}[Interval soundness for every network]\label{thm:box_sound}
\lean{OpenTheory.Exact.Net.box_sound, OpenTheory.Exact.reluMLP_box_sound}
\leanok
\uses{def:interval_rule, lem:exact_primitives}
For every network whose primitives carry interval rules, valid parameters, and every real input $x$
in a box $b$, the TorchLean model sends $x$ into the box computed by interval propagation from $b$.
In particular this holds for every ReLU multilayer perceptron with integer weights, of any depth
and widths.
\end{theorem}
\begin{proof}
\leanok
Induction on the composition, as in Theorem~\ref{thm:exact_agrees}.
\end{proof}

\begin{theorem}[Region certificate and local robustness]\label{thm:region_certificate}
\lean{OpenTheory.Predicts, OpenTheory.OnRegion.certify, OpenTheory.OnRegion.certify_sound, OpenTheory.OnRegion.robust_of_certify, OpenTheory.Exact.inBox_lower, OpenTheory.OnRegion.inBox_ball}
\leanok
\uses{thm:box_sound}
Say that scores $y$ predict class $c$ if $y_j<y_c$ for every $j\ne c$. The integer check
$\mathrm{certify}$ accepts when the parameters are valid, the input box is well formed
($\ell_i\le u_i$, one pair per coordinate) and, in the propagated output box, the lower bound of
class $c$ exceeds the upper bound of every other class. If it accepts, the box is not empty and the
TorchLean model predicts $c$ on every real input in it. For the box of radius $\varepsilon$ around an
integer input $x_0$, this reads: every real $x$ with $|x_i-x_{0,i}|\le\varepsilon$ for all $i$ is
classified as $c$.
\end{theorem}
\begin{proof}
\leanok
Interval soundness puts the output in the propagated box; there the score of $c$ is at least its
lower bound, which exceeds every other upper bound. The lower corner of a well-formed box lies in
it.
\end{proof}

\bibliographystyle{alpha}
\bibliography{references}
